%% 第31回編集委員会（物理第5巻）：組 原W4
%% p5-07.tex  原-7　電子の席の決まり：原子と周期表（★2〜5・3段）
%% 第32回（修5-3、H-3）：目でわかる ①② ★2 → 手でわかる ③〜⑨ ★3 → 式でわかる ⑩〜⑫ ★5。③の P(r) の山を微分で求める所と、
%%   積分で確率・平均を求める所は「もっと知りたい★4」へ移した。直しの記録は tools/FIX-phys5-review-3.md。
%% 中身は計画書 tools/phys5-plan.md の B「原-7」。公式の持ち主は A（A5・A41・A52・A53・A86）。
%%   ほかの組の章は、計画書 B の節の番号で指した（原-4 ⑤⑥⑩、原-5 ③④⑥⑦、原-6 ⑧⑨⑩⑪、原-8 ④⑤⑦⑫、原-9、原-11）。
%%   波動関数の書き方は計画書 C-1 のとおり（$\psi$、時間をふくむときは $\Psi$。虚数単位はこの章では使わない）。
%% 数値はすべて Python で計算した：checks/物理第5巻/原-7/calc_p5_07.py（出力 calc_p5_07_out.txt）、
%%   code_p5_07.py（⑩の囲みのプログラム）、kanji_p5_07_08.py（常用漢字の確かめ）。電子の雲の点は cloud.py で作った（同じフォルダ）。
%% 前の巻・数学の本との約束（計画書 A-6・C-6-2）
%%   ・Q12（熱-15 ⑦「磁場の中の電子・原子核の向きが2つしかない（なぜかは『原子』で）」）→ ⑥・例題6（熱-15 の「上の数/下の数」の形）
%%   ・Q15（熱-15 ⑨「電子は1つの状態に1個まで（原-7）」）→ ⑦（「1つの状態（棚）」の言い方を受けた）
%%   ・Q53（電-6 の MRI 1.5〜3 T）→ ⑥・例題7（63.9 MHz・127.7 MHz）
%%   ・Q54（電-7 ⑩ 交換相互作用・磁区のしくみ）→ ⑧（交換相互作用のもと）。磁区そのものと帯の磁性は原-8 ⑫
%%   ・Q41（電-1 の「原子の間の力は量子力学で（原-7）」）→ ⑪
%%   ・Q4（力-8 の原子のばね）→ ⑪（結合のエネルギーの谷。ばね定数そのものは原-6 ⑨）
%%   ・Q10（熱-14 の等分配：分子の回転）→ ⑪・例題11（振動は原-6 ⑩）
%%   ・M1（数学 4-15:698「電子の雲は球座標で簡単な式。周期表（原-7）」）→ ②⑩
%%   ・M17（数学 8-12⑦ S6 の球のラプラシアン）→ ⑩
%%   ・電-4 の銅の自由電子 n＝8.49×10²⁸ /m³（計画書 C-4 の2）→ 例題8（同じ数になる）
%% 記号（この章）：主量子数 n、方位量子数 l、磁気量子数 m_l、スピン量子数 m_s、原子番号 Z、有効核電荷 Z_eff、
%%   ボーア半径 a_0、クーロンの定数 k（電-1 と同じ。この章では波数を使わない）、ボルツマン定数 k_B（かならず添字）、
%%   エネルギー E（電離エネルギーは E_i とことば）、位置エネルギー U（電圧 V と分ける）、波動関数 ψ、
%%   ボーア磁子 μ_B、g 因子 g（重力加速度はこの章に出ない）、磁場 B（磁束密度。電-6 と同じ）、磁化 M（電-7 と同じ）、
%%   振動数（電波）は「周波数」f（計画書 E の5）、換算質量 μ（⑪・例題11 だけ。μ_B とは添字で分ける）、慣性モーメント I（電流は記号にしない）。
%% 出典（11月11日に確かめた）
%%   ・定数（h・ħ・e・c・k_B・N_A・m_e・ε_0・a_0・R_∞hc＝13.605693 eV・μ_B・電子の g 因子 2.00231930436・陽子の磁気回転比/2π＝42.577478461 MHz/T・u・μ_0）：
%%       CODATA 2022 https://physics.nist.gov/cuu/Constants/Table/allascii.txt（計画書 D4）
%%   ・電離エネルギー（H〜Xe の中性原子）：NIST Atomic Spectra Database, Ionization Energies
%%       https://physics.nist.gov/cgi-bin/ASD/ie.pl?spectra=H-Xe+I&units=1 （curl で表を読んだ。He 24.587、Li 5.392、Ne 21.565、Na 5.139、Ar 15.760、K 4.341 eV）
%%   ・CO の分子の定数（B_e＝1.93128 cm⁻¹、α_e＝0.01750 cm⁻¹、r_e＝1.128323 Å、ω_e＝2169.8 cm⁻¹）、H₂ の D_0＝36118 cm⁻¹（4.478 eV）：
%%       NIST Chemistry WebBook https://webbook.nist.gov/cgi/cbook.cgi?ID=C630080&Mask=1000 、ID=C1333740&Mask=1000（curl で読んだ）
%%   ・密度（銅 8.960、鉄 7.874 g/cm³）：NIST X-Ray Mass Attenuation Coefficients 表1 https://physics.nist.gov/PhysRefData/XrayMassCoef/tab1.html （計画書 D7）
%%   ・ノーベル賞：1945 パウリ（排他律）、1943 シュテルン（分子線の方法・陽子の磁気モーメント）、1952 ブロッホとパーセル（核磁気の精密な測定）、
%%       1922 ボーア：https://www.nobelprize.org/prizes/physics/年/summary/（curl で読んだ）
%%   ・113番元素（ニホニウム、Nh）：理化学研究所の森田浩介らのグループ。2004年7月・2005年4月・2012年8月に合成、2015年12月31日に IUPAC・IUPAP が発見を認定、
%%       2016年11月30日に IUPAC が名前を決定。理研（www.riken.jp）は作業環境から届かなかった。検索の抜粋2件（理研のプレスリリース 20161130_1、
%%       JST サイエンスポータル 20161201_02）で照らし合わせた。
%%   ・118番までで第7周期がうまったこと、超重元素の研究の現状：K. Hagino, "Superheavy Elements: Beyond the 7th Period in the Periodic Table", arXiv:1812.05805（要旨を curl で読んだ）。
%%       「安定の島」（Z＝114・N＝184 のあたり）にはまだ届いていない（いちばん中性子の多いフレロビウムでも 9 個足りない：同じ論文の本文の抜粋）。
%%   ・メンデレーエフの周期表 1869年：2019年が「周期表の国際年」（150周年。国連総会が 2017年に宣言）であることを、BIPM の記事（www.bipm.org/en/-/2019-mendeleev の検索の抜粋）
%%       と IUPAC の行事の資料の抜粋で照らし合わせた。
%% 〔要確認〕（確かめられなかった値）
%%   ・鉄の1原子あたりの磁気モーメント 2.2 μ_B（固体物理の教科書 Kittel の値）：一次資料に届かなかった。本文では「はかると約 2.2」と書き、
%%       電-7 の飽和 約 2 T（日本製鉄の電磁鋼板の資料）とつき合わせた（例題9）。
%%       出典の組（11月11日）：今回届くようになった一次資料に固体の磁気モーメントの表はなく、論文（doi.org の先の出版社）は届かなかった。〔要確認〕のまま。
%%   ・H₂⁺ の結合エネルギー 2.65 eV は本文に数を出さない（「H₂ の半分より少し大きい」とだけ書いた）。
\makeatletter
% 電子の席（箱とスピンの矢印）  \txhbox{(左下)}{上向きの数 0,1}{下向きの数 0,1}
\newcommand\txhbox[3]{\begin{scope}[shift={#1}]
  \filldraw[fill=white, draw=bkInk, line width=0.6pt] (0,0) rectangle (6,6);
  \ifnum#2>0 \draw[bkBlue!80!black, line width=1.0pt, -{Stealth[length=1.6mm]}] (2,1) -- (2,5.2);\fi
  \ifnum#3>0 \draw[bkRed!80!black, line width=1.0pt, -{Stealth[length=1.6mm]}] (4,5) -- (4,0.8);\fi
  \end{scope}}
% 周期表のます  \txhcell{列}{行}{色}{文字}
\newcommand\txhcell[4]{\filldraw[fill=#3, draw=bkInk!70, line width=0.35pt] ({(#1-1)*6.3},{-(#2-1)*5.6}) rectangle ++(6.3,-5.6);
  \node[font=\sffamily\footnotesize, text=bkInk, inner sep=0pt] at ({(#1-1)*6.3+3.15},{-(#2-1)*5.6-2.8}) {#4};}
\tikzset{txh lab/.style={font=\sffamily\footnotesize, text=bkInk}, txh arr/.style={line width=0.9pt, -{Stealth[length=2.2mm]}},
  txh dim/.style={bkGray, line width=0.5pt, {Stealth[length=1.8mm]}-{Stealth[length=1.8mm]}}}
\makeatother

\phchapter{7}{電子の席の決まり：\\原子と周期表}{ネオンは、なぜほかの物質とほとんど結びつかない？}{\bklevelrange{2}{5}}{電子の席の決まり：原子と周期表}
\bkmichishirube{3-17,4-15,7-8,8-12,熱-15,電-7,原-4,原-5,原-6}{原-8,原-9,原-11,原-13}

\begin{goal}
\begin{itemize}[label={\color{bkTheme}●}, leftmargin=1.6\zw, itemsep=0pt, topsep=1pt]
\item \bklevel{2}\enspace 周期表の横の並びの数（2・8・8・18…）が、原子の中の電子の「席」の数で決まることを説明できる。
\item \bklevel{3}\enspace 量子数で席を数えて $2n^2$ を導き、ナトリウム・鉄・銅などの電子配置を書ける。
\item \bklevel{4}\enspace 水素原子の電子の雲から、電子が見つかりやすい距離や見つかる確率を求め、電離エネルギーから有効核電荷を見積もれる。
\item \bklevel{3}\enspace 磁場の中の電子や陽子の向きが2つに分かれることから、ESR や MRI の周波数を求められる。
\item \bklevel{5}\enspace 3次元のシュレディンガー方程式に水素原子の $1s$ の波動関数を代入して $-13.6$ eV を確かめ、化学結合や分子の回転のしくみを読める。
\end{itemize}
\end{goal}

\phkoushiki{%
\phk{電子殻に入る電子の数は、内側から 2・8・18（$2n^2$）。電子配置と周期表}{高校化学基礎}{ph:原-7:2n2}
\phk{原子が電子を失うと陽イオン、受け取ると陰イオンになる}{中学理科}{ph:原-7:ion}
}

\begin{hatena}
\noindent\begin{minipage}[c]{0.56\linewidth}\raggedright
ネオンサインの赤い光は、ガラスの管に入れたネオンという気体が出す光です。ネオンは、空気の中にもほんの少しふくまれていますが、ほかのどんな物質とも、ほとんど結びつきません。水素は酸素と結びついて水になり、ナトリウムは塩素と結びついて食塩になるのに、です。

ネオンの原子は電子を10個、ナトリウムの原子は11個もっています。たった1個しか違わないのに、なぜ、こんなに性質が違うのでしょう。
\end{minipage}\hfill
\begin{minipage}[c]{0.41\linewidth}\centering
\begin{tikzpicture}[x=0.92mm, y=0.92mm]
  \foreach \cx/\nm/\nn in {10/ネオン/10, 44/ナトリウム/11} {
    \fill[bkRed!70!black] (\cx,14) circle (1.6);
    \draw[bkInk, line width=0.5pt] (\cx,14) circle (4.2) (\cx,14) circle (8.4);
    \node[txh lab, anchor=north, align=center] at (\cx,1.2) {\nm\\（\nn 個）};}
  \draw[bkInk, line width=0.5pt] (44,14) circle (12.0);
  \foreach \a in {90,270} {\fill[bkBlue!80!black] ($(10,14)+(\a:4.2)$) circle (0.9); \fill[bkBlue!80!black] ($(44,14)+(\a:4.2)$) circle (0.9);}
  \foreach \a in {0,45,...,315} {\fill[bkBlue!80!black] ($(10,14)+(\a:8.4)$) circle (0.9); \fill[bkBlue!80!black] ($(44,14)+(\a:8.4)$) circle (0.9);}
  \fill[bkBlue!80!black] ($(44,14)+(60:12)$) circle (0.9);
  \node[txh lab, anchor=south] at (44,26.5) {1個だけ外に};
\end{tikzpicture}
\end{minipage}
\end{hatena}

\noindent 第2部の前の2章（原-5・原-6）では、1次元の箱や井戸に閉じこめた電子が、とびとびのエネルギーの「定在波」になることを見てきました。この章では、いよいよ本物の原子の中の電子を調べます。原子の中の電子は、原子核の電気の力で閉じこめられた波です。その波には決まった形がいくつかあり、それぞれの形には電子が1個ずつしか入れません。この「席」の決まりから、元素の性質が周期的に繰り返されること、つまり\textbf{周期表}の形が出てきます。\bkidx{しゅうきひょう}{周期表}化学で習う電子殻の 2・8・18 も、MRI が体の中をうつせるわけも、磁石が磁石であるわけも、もとは同じ決まりです。

\begin{bkfigure}
\begin{tikzpicture}[x=1mm, y=1mm]
  % 箱（原-5）
  \begin{scope}[shift={(0,0)}]
    \draw[bkInk, line width=1.0pt] (0,30) -- (0,0) -- (30,0) -- (30,30);
    \foreach \yy/\nn in {2.6/1, 10.4/2, 23.4/3} {\draw[bkBlue!80!black, line width=0.9pt] (0,\yy) -- (30,\yy); \node[txh lab, anchor=west] at (31,\yy) {$n=\nn$};}
    \node[txh lab, anchor=north, align=center] at (15,-1.5) {箱（原-5）\\1次元};
  \end{scope}
  % 有限の井戸（原-6）
  \begin{scope}[shift={(54,0)}]
    \draw[bkInk, line width=1.0pt] (-4,24) -- (4,24) -- (4,0) -- (26,0) -- (26,24) -- (34,24);
    \foreach \yy in {4.2, 15.6} {\draw[bkBlue!80!black, line width=0.9pt] (4,\yy) -- (26,\yy);}
    \node[txh lab, anchor=north, align=center] at (15,-1.5) {有限の井戸（原-6）\\1次元};
  \end{scope}
  % 原子核のすりばち（原-7）
  \begin{scope}[shift={(112,0)}]
    \begin{scope}
      \clip (-19,-0.5) rectangle (19,31);
      \draw[bkInk, line width=1.0pt, domain=-19:-0.15, samples=60, smooth] plot (\x,{28*(1-1/(1+abs(\x)/4))});
      \draw[bkInk, line width=1.0pt, domain=0.15:19, samples=60, smooth] plot (\x,{28*(1-1/(1+abs(\x)/4))});
      \foreach \yy/\ww in {6.0/1.12, 22.5/16.4, 25.56/19} {\draw[bkBlue!80!black, line width=0.9pt] (-\ww,\yy) -- (\ww,\yy);}
    \end{scope}
    \fill[bkRed!70!black] (0,0) circle (1.0);
    \node[txh lab, anchor=north, align=center] at (0,-1.5) {原子核のすりばち（原-7）\\3次元、$U=-ke^2/r$};
  \end{scope}
\end{tikzpicture}
\bkfigcaption{第2部で電子を閉じこめた井戸（太い線が位置エネルギーの形、細い横線がとびとびのエネルギー）。この章では、原子核の電気の力が作る、深さが距離に反比例する「すりばち」に、3次元で電子を閉じこめる。段の間かくは、上へ行くほどつまる。}\label{ph:原-7:fig-wells}
\end{bkfigure}

\begin{yattemiyou}[元素のカードを、似たものどうしで並べよう]
\textbf{用意するもの}\par
厚紙で作ったカード20枚（1枚に1つの元素。原子番号・名前・常温のようす・ほかの物との結びつきやすさを、資料で調べて書く）、はさみ、筆記用具\par\vspace{2pt}
\begin{enumerate}[label={\bkcircled[bkGreen]{\arabic*}}, leftmargin=2.4\zw, itemsep=1pt, topsep=0pt]
\item 原子番号1（水素）から20（カルシウム）まで、カードを原子番号の順に1列に並べる。
\item 「ほとんど結びつかない気体」のカード（ヘリウム・ネオン・アルゴン）を見つけ、そのすぐあとで列を折り返して、次の行に進む。
\item できた表で、たてに並んだカードどうしの性質が似ているかを見る。
\end{enumerate}

\par\vspace{2pt}\textbf{気をつけること}：ナトリウムやカリウムを、自分で水に入れてはいけません（激しく反応して危険です）。性質は、資料や先生の実験で調べましょう。
\end{yattemiyou}

\begin{bkcenter}
\begin{bktable}{colspec={Q[c,m,wd=13mm] Q[l,m,wd=32mm] Q[l,m,wd=58mm] Q[c,m,wd=14mm]}, rowsep=0.5pt, row{1}={bg=bkGreen!10}, cells={font=\small}}
番号 & 名前 & 性質（いくつか） & 行の長さ \\
1・2 & 水素・ヘリウム & ヘリウムは結びつかない気体 & 2 \\
3〜10 & リチウム〜ネオン & リチウムは水と激しく反応。フッ素は激しく結びつく。ネオンは結びつかない & 8 \\
11〜18 & ナトリウム〜アルゴン & ナトリウムはリチウムに似る。塩素はフッ素に似る。アルゴンは結びつかない & 8 \\
19・20 & カリウム・カルシウム & カリウムはナトリウムに似る（水と激しく反応） & ― \\
\end{bktable}
\bktabcaption{やってみようで、ソラさんが資料で調べて作った表}\label{ph:原-7:tab-yattemiyou}
\end{bkcenter}

\noindent やってみようの答え：表\ref{ph:原-7:tab-yattemiyou}のように折り返すと、行の長さは 2・8・8 になり、たてに並んだ元素の性質がよく似ます（リチウム・ナトリウム・カリウム、フッ素・塩素、ヘリウム・ネオン・アルゴン）。1869年に、ロシアのメンデレーエフは、当時知られていた元素を同じように並べ、性質の周期的な繰り返しを表にまとめました（2019年は、その150年目として国連の「周期表の国際年」になりました）。でも、\textbf{なぜ}2・8・8なのかは、それから約60年後、量子力学ができるまでわかりませんでした。この章で、その「なぜ」に答えます。

\begin{omoidasou}
\bkrecall{原-4}{水素原子の電子のエネルギーは $E_n=-\dfrac{13.6}{n^2}$ eV（$n=1,2,3,\ldots$）。とびとびのエネルギーの差が、出す光のエネルギー $hf$ になる。ボーアの円の半径は $r_n=n^2a_0$、$a_0=\qty{0.0529}{nm}$（ボーア半径）。}
\bkrecall{原-5}{電子は波としてふるまい、波動関数 $\psi$ で表す。電子が見つかる確率は $|\psi|^2$ に比例する。閉じこめた電子は定在波になり、エネルギーはとびとびになる。時間によらないシュレディンガー方程式 $-\dfrac{\hbar^2}{2m}\dfrac{\dd^2\psi}{\dd x^2}+U\psi=E\psi$。}
\bkrecall{熱-15}{エネルギーの差が $\Delta\varepsilon$ の2つの準位では、上の数と下の数の比は $e^{-\Delta\varepsilon/(k_{\mathrm B}T)}$。300 K の $k_{\mathrm B}T$ は約 0.0259 eV。}
\bkrecall{電-6}{⑪：磁場 $B$ の中のコイル（磁石の強さ $NIS$）は、トルク $NISB\sin\theta$ を受け、面に垂直な向きが磁場の向きにそろうように回る。そろった向きが、エネルギーの低い向き。}
\bkrecall{4-15}{球座標 $(r,\theta,\varphi)$：$r$ は原点からの距離、$\theta$ はたての角、$\varphi$ は横の角。}
\bkrecall{3-17}{確率は「起こる場合÷全部の場合」。たくさん繰り返すと、起こった割合が確率に近づく。}
\end{omoidasou}

\noindent この章の記号：原子番号は $Z$、量子数は $n$・$l$・$m_l$・$m_s$ と書きます。質量の $m$ と区別するため、電子の質量は $m_{\mathrm e}$ と添字をつけます。電波や磁石のゆれの「1秒に繰り返す回数」は、電気の本と同じく\textbf{周波数} $f$ とよびます。

%% ======================== 物理の道具 ========================
\bksection{doguu}{物理の道具}

\phstage{me}{周期表と電子の雲を見る}{2}

\bksubsection{1}{周期表の形を見る}
\label{ph:原-7:sec1}
いまの周期表には、118の元素が並んでいます。図\ref{ph:原-7:fig-pt}は、その形だけを取り出したものです。横の行を\textbf{周期}、たての列を\textbf{族}といいます。\bkidx{しゅうき}{周期（周期表）}行の長さは、上から 2・8・8・18・18・32・32 です。どの行も、右のはしは、ほとんど結びつかない気体（ヘリウム・ネオン・アルゴン・クリプトン・キセノン・ラドン）で終わります。これらをまとめて\textbf{貴ガス}といいます。\bkidx{きがす}{貴ガス}

\begin{bkfigure}
\begin{tikzpicture}[x=1mm, y=1mm]
  % s ブロック
  \foreach \r in {1,...,7} {\txhcell{1}{\r}{bkBlue!14}{}}
  \foreach \r in {2,...,7} {\txhcell{2}{\r}{bkBlue!14}{}}
  % d ブロック
  \foreach \r in {4,...,7} {\foreach \cc in {3,...,12} {\txhcell{\cc}{\r}{bkGreen!14}{}}}
  % p ブロック
  \foreach \r in {2,...,7} {\foreach \cc in {13,...,18} {\txhcell{\cc}{\r}{bkOrange!16}{}}}
  \txhcell{18}{1}{bkBlue!14}{}
  % f ブロック（下に別に）
  \foreach \r in {9,10} {\foreach \cc in {4,...,17} {\txhcell{\cc}{\r}{bkPurple!12}{}}}
  % 元素の記号（いくつか）
  \node[txh lab] at (3.15,-2.8) {H};   \node[txh lab] at (110.25,-2.8) {He};
  \node[txh lab] at (3.15,-8.4) {Li};  \node[txh lab] at (110.25,-8.4) {Ne};
  \node[txh lab] at (3.15,-14.0) {Na}; \node[txh lab] at (110.25,-14.0) {Ar};
  \node[txh lab] at (3.15,-19.6) {K};  \node[txh lab] at (47.25,-19.6) {Fe}; \node[txh lab] at (66.15,-19.6) {Cu}; \node[txh lab] at (110.25,-19.6) {Kr};
  \node[txh lab] at (110.25,-25.2) {Xe}; \node[txh lab] at (110.25,-30.8) {Rn}; \node[txh lab] at (110.25,-36.4) {Og};
  \node[txh lab] at (78.75,-36.4) {Nh};
  % ブロックの名前
  \node[font=\sffamily\bfseries\small, text=bkBlue!70!black] at (6.3,-42.0) {s};
  \node[font=\sffamily\bfseries\small, text=bkGreen!50!black] at (47.25,-14.0) {d ブロック};
  \node[font=\sffamily\bfseries\small, text=bkOrange!70!black] at (91.35,-2.8) {p ブロック};
  \node[font=\sffamily\bfseries\small, text=bkPurple!70!black] at (66.15,-60.0) {f ブロック（第6・7周期の、3族の位置に入る）};
  % 貴ガスの列を太線で
  \draw[bkInk, line width=1.4pt] (107.1,0) rectangle (113.4,-39.2);
  % 行の長さ
  \foreach \r/\L in {1/2,2/8,3/8,4/18,5/18,6/32,7/32} {\node[txh lab, anchor=west] at (115,{-(\r-1)*5.6-2.8}) {\L};}
  \node[txh lab, anchor=south west, align=left] at (114,1.0) {行の長さ};
  \node[txh lab, anchor=east] at (106,-42.0) {太線の列：貴ガス};
\end{tikzpicture}
\bkfigcaption{周期表の形。左はしの2列が s ブロック、まん中の10列が d ブロック、右の6列が p ブロック、下に別に置いた2行（14列）が f ブロック。ヘリウム（He）は s ブロックだが、貴ガスの列に置く。Nh はニホニウム（113番）、Og はオガネソン（118番）。}\label{ph:原-7:fig-pt}
\end{bkfigure}

周期表には、もう1つ大事な形があります。左から2列、まん中の10列、右の6列、そして下に別に置いた14列という、\textbf{ブロック}の分かれ方です（図\ref{ph:原-7:fig-pt}の s・d・p・f）。2・6・10・14 という数は、どこから来るのでしょう。2・8・18・32 は、それぞれ
\[ 2=2,\qquad 8=2+6,\qquad 18=2+6+10,\qquad 32=2+6+10+14 \]
と、ブロックのはばの和になっています。そして、$2=2\times1^2$、$8=2\times2^2$、$18=2\times3^2$、$32=2\times4^2$ です。この「2・6・10・14」と「$2n^2$」が、電子の席の数です。④で、その数え方を作ります。

\begin{phmijika}[ネオンサインと、電球の中のアルゴン]
貴ガスは、ほかの物と結びつかないことが役に立ちます。白熱電球の中には、アルゴンなどを入れて、熱いタングステンのフィラメントが酸素と結びついて燃えつきないようにしています。溶接では、とけた金属が空気とふれないように、アルゴンやヘリウムを吹きつけます。ネオンサインは、管の中のネオンに電圧をかけて、決まった色の光（原-4 のスペクトル）を出させます。
\end{phmijika}

\Needspace{12\baselineskip}\bksubsection{2}{殻と、電子の雲の形}
\label{ph:原-7:sec2}
化学の教科書では、原子核の周りに同じ中心の円をかき、内側から K 殻・L 殻・M 殻……とよび、電子を「殻」に入れて数えます（はてな？の図）。殻の絵は、席の数を数えるにはとても便利です。でも、電子は、円の上を回る小さな粒ではありません。原-5 で見たように、電子は波で、どこに見つかるかは確率で決まります。

図\ref{ph:原-7:fig-cloud}は、水素原子の電子が見つかる場所を、何度もはかったとして点で打ったものです（点の濃いところほど、見つかりやすい）。いちばん低いエネルギーの状態（$1s$）は、原子核を中心にした、ふちのぼやけた丸い雲です。その次のエネルギーの状態には、丸い雲の中に「電子が見つからない球面」（節）がある $2s$ と、まん中で2つに分かれた、ひょうたんのような $2p$ があります。雲は3次元の形です。図は、原子核を通るうすい板の中の点だけを打ちました。

\begin{bkfigure}
\begin{tikzpicture}[x=1mm, y=1mm]
  \begin{scope}[shift={(18,18)}]
    \foreach \px/\py in {0.4/-0.0,0.6/0.7,5.1/-2.7,8.5/1.5,1.4/1.2,-0.5/1.8,1.1/-4.7,-4.2/-2.3,-3.6/0.3,0.3/2.3,-6.8/-6.4,3.2/-0.9,-1.8/-8.0,-3.1/-5.2,-2.0/-0.3,1.3/0.4,0.3/-2.2,0.7/2.7,-1.6/0.5,-9.2/-5.3,-3.0/-2.5,1.8/-1.6,-3.2/0.0,-1.2/1.5,-0.4/0.1,1.4/-0.5,-2.5/-2.7,2.6/-0.7,-2.8/6.4,8.2/-4.8,-5.1/-11.6,-4.4/-0.8,0.7/1.5,1.3/6.1,1.3/-2.4,0.2/0.7,2.8/5.9,4.5/-9.3,-1.1/2.8,-5.5/-5.6,1.0/-1.6,-0.8/1.7,1.6/3.3,-2.6/-0.7,3.1/3.2,-2.3/-0.6,0.7/-0.1,0.9/1.7,15.9/-7.3,0.4/-2.7,-2.5/-0.9,-8.6/-2.1,-5.6/0.9,-2.9/0.8,1.9/2.0,-5.7/11.0,4.1/0.0,-7.2/6.2,4.1/-2.4,3.1/2.1,-3.4/-2.7,4.4/-0.5,-4.1/-2.8,-1.1/-2.5,0.2/-1.6,3.1/2.1,-0.3/-1.4,-0.1/-0.5,-0.1/-0.9,7.4/11.3,5.8/-4.8,0.8/-3.1,-2.2/1.9,-1.3/-0.3,-2.4/-0.4,-2.7/-8.8,3.1/2.2,0.9/2.6,1.4/1.0,5.2/4.1,-2.2/-1.0,0.2/-4.5,3.0/0.2,-1.5/4.5,-4.2/0.5,6.5/-5.8,-2.6/-0.1,0.9/-3.2,-6.2/-2.3,3.8/-4.4,-1.1/0.4,3.2/0.9,6.0/-1.3,-4.0/-4.5,1.7/-2.6,1.8/-2.2,3.8/-0.9,3.9/-1.3,-4.4/-1.2,0.6/1.1,2.6/1.6,-5.7/6.1,-0.7/-2.9,0.5/-0.3,-1.8/-2.7,2.1/-0.2,-2.2/2.3,-0.3/4.9,-0.1/-1.5,-0.2/2.0,0.2/0.9,1.2/-0.9,-3.5/-3.5,-0.2/0.7,3.6/1.0,-2.4/0.4,-6.2/-5.1,3.3/-5.0,1.1/-0.7,1.1/-0.8,-3.8/0.4,0.2/5.0,0.5/-1.6,2.3/-0.4,1.8/3.7,-6.6/-2.3,2.8/3.9,0.1/1.8,-0.1/0.6,-0.2/2.5,7.6/-1.5,2.4/1.0,-2.7/2.4,2.2/1.8,-0.5/-2.2,5.4/-1.4,4.8/-0.3,1.0/-0.6,-0.4/2.6,-8.7/-0.9,10.0/2.5,3.6/-8.8,-0.4/4.3,0.1/1.5,-1.7/-1.3,-0.6/0.5,-2.9/-6.2,-3.0/3.2,3.9/-1.0,5.4/-6.6,-3.8/-8.7,1.8/7.6,2.4/6.0,-6.4/0.3,0.8/-1.9,5.3/4.9,-0.6/-0.3,-1.8/1.5,-5.3/4.3,6.2/-2.2,-2.1/2.5,-3.8/2.6,-6.1/-1.3,-2.9/5.2,-2.8/0.7,6.8/5.4,-0.6/2.0,4.5/5.6,-7.0/-8.9,-1.7/-0.0,0.7/3.2,-3.4/-0.6,-3.3/2.0,-0.3/3.6,2.5/-2.8,2.0/-2.2,-0.2/-7.0,-2.6/1.7,3.1/6.5,-0.4/3.8,-1.6/-0.8,-3.3/0.7,-6.0/-4.5,0.7/1.7,1.0/7.5,-0.1/-0.8,0.3/-1.5,-3.1/1.3,-8.2/-4.6,-5.8/9.8,1.5/-2.6,-7.1/5.8,4.9/0.8,-4.2/-4.9,0.8/-4.0,10.9/-0.1,-4.2/-7.8,-11.3/-4.1,-2.2/-13.4,1.5/6.4,-2.5/0.2,4.1/0.5,-3.2/-2.9,-1.6/-1.4,0.5/-0.1,-1.3/1.2,2.3/1.3,0.8/-0.7,-4.4/3.6,-0.4/1.7,-3.0/0.1,-1.6/4.9,7.5/2.6,-3.6/-1.4,-0.8/0.2,2.6/7.0,-3.2/2.9,0.2/2.0,-2.5/-3.6,-0.2/0.8,0.9/-3.5,3.6/0.6,-0.5/-0.3,0.3/0.7,0.1/-2.9,3.7/-0.6,2.5/-0.7,-1.1/-0.0,-2.1/0.2,-1.5/-1.5,2.2/1.2,7.1/4.2,1.4/-8.1,8.9/10.5,0.4/-0.9,4.0/3.3,-0.8/-3.1,-8.7/-4.1,-0.6/-1.1,1.2/-1.5,2.9/10.2,4.1/-5.6,1.2/4.5,0.5/0.1,-0.6/3.8,-7.1/6.3,-0.1/-0.2,-3.8/1.2,3.9/-0.3,4.4/0.6,5.9/-6.6,-0.1/0.8,-3.1/5.9,2.3/-1.0,-3.6/3.7,-0.5/3.2,2.3/-0.1,-7.6/-2.7,-2.2/-0.4,-1.4/-0.2,1.3/-0.7,1.9/-2.8,2.2/1.0,4.6/-0.1,-6.2/3.8,0.4/-1.5,-3.8/-0.1,1.9/-1.0,-0.0/0.8,-2.2/3.5,-2.1/2.1,-3.4/2.9,1.7/-6.5,-0.8/1.3,-3.6/-3.6,0.9/1.1,4.2/-3.3,0.4/3.8,-3.0/5.2,-1.2/1.6,1.9/0.2,3.0/-8.7,-2.6/-4.3,1.7/2.9,-0.3/2.3,-4.0/1.9,3.5/-7.1,1.4/-0.7,1.5/-3.1,1.6/-0.5,-5.2/-0.8,6.2/-4.9,1.1/0.5,-3.3/-0.1,-1.8/1.3,-0.8/4.3,3.4/-1.1,3.4/5.1,-3.3/1.7,1.2/-1.7,-1.8/-0.4,7.3/-3.5,-5.1/0.9,-4.6/6.8,-0.1/-1.6,-4.4/-6.3,2.3/-1.0,-5.4/1.3,-3.6/-3.9,-2.4/2.9,2.5/-2.5,5.6/-2.0,-0.5/1.9,-0.2/2.5,-4.7/-10.0,-0.8/0.2,0.1/-1.7,1.7/-0.3,-0.3/1.6,-1.5/-1.9,0.5/-0.8,-8.6/-3.0,1.1/-2.3,-1.4/-1.5,3.1/-3.8,0.3/0.5,-0.7/-0.3,6.6/3.0,3.4/-0.6,6.9/0.3,0.5/0.2,2.5/4.2,2.3/-6.7,0.3/2.3,2.8/-2.2,2.5/6.6,-3.6/0.9,1.7/-3.2,7.6/-3.3,1.0/-1.9,9.6/9.8,2.0/3.6,-0.6/-1.8,3.6/0.1,1.5/-0.4,-2.1/-1.5,-4.0/1.2,1.0/0.0,-3.8/-2.4,-1.0/-0.4,-0.6/0.2,-1.0/-1.6,-0.0/1.2,-5.5/0.3,-3.5/0.3,-3.6/0.1,1.6/3.1,3.4/-1.3,4.4/-4.2,0.3/-0.1,-3.2/3.9,2.2/3.6,1.4/-2.5,0.0/3.5,-3.9/2.3,-0.7/1.2,-5.9/0.8,2.2/0.9,-2.4/-1.6,1.4/2.0,4.5/-0.8,-0.2/0.2,1.4/0.1,-3.9/1.5,7.7/0.1,1.4/1.4,0.6/1.5,-5.9/1.7,0.0/0.0,-1.1/0.1,1.3/-8.0,2.5/4.1,-6.8/0.6,-2.5/6.2,5.5/0.5,-0.2/2.8,-5.4/1.6,-4.4/2.6,1.9/11.3,-0.2/6.6,1.9/0.1,2.2/2.6,-0.6/0.6,2.6/2.9,1.6/-2.6,-2.6/2.2,1.3/-1.8,-2.0/4.1,-4.7/-2.5,-1.2/3.8,4.6/-0.6,2.8/-1.1,3.6/-3.4,-2.0/-6.8,-1.7/-2.4,0.2/-5.8,0.2/-6.0,-0.5/0.4,1.2/-0.0,3.9/1.7,7.7/3.0,-1.5/-2.3,1.6/3.2,4.9/-9.1,1.6/-0.6,1.7/0.8,0.7/-3.2,-0.8/1.1,-10.5/-4.4,3.3/4.7,1.9/-8.6,-7.5/-1.7,-3.4/1.6,-0.2/6.5,-4.1/-2.4,1.4/1.3,-2.4/-1.6,-6.3/-5.1,9.3/-3.7,-11.4/0.8,0.7/-4.8,-0.2/-3.0,3.6/-2.5,1.5/-1.0,3.4/-4.8,-4.0/-3.8,3.0/-10.0,3.2/-3.7,9.2/2.4,-3.0/1.7,3.6/-0.3,-0.9/-0.1,-1.4/-5.1,-1.8/-3.3,0.5/-0.7,0.6/-4.9,-1.9/4.2,-0.3/-0.4,-0.6/0.9,-2.7/-1.1,0.4/-8.2,-0.1/-2.9,1.4/-0.6,1.5/-0.4,-5.2/1.9,-2.2/4.2,-0.2/-3.5,-5.1/8.4,3.6/-5.6,-9.4/-1.5,-0.4/-0.3,0.9/-3.0,-2.2/-3.9,0.8/-0.8,0.7/2.3,-3.8/-0.3,-2.2/2.5,0.9/1.8,0.3/-2.6,-0.2/2.8,1.7/-1.1,3.5/0.1,1.6/-0.6,-3.3/2.0,-5.0/3.0,0.1/2.0,-5.6/-3.1,7.7/1.3,1.3/-0.6,-1.0/1.6,-2.0/1.6,-2.7/-7.6,-5.3/3.6,-4.6/-2.4,-0.6/5.3,1.1/1.2,-4.6/3.3,-7.0/-11.1,-12.0/2.0,6.9/-7.4,2.2/3.8,1.2/-4.3,0.2/6.0,-4.2/-7.1,-2.6/2.0,3.0/-3.7,0.5/-4.1,4.3/4.7,1.0/1.6,2.6/3.3} {\fill[bkBlue!80!black] (\px,\py) circle (0.32);}
    \fill[bkRed!80!black] (0,0) circle (0.8);
    \node[txh lab, anchor=south] at (0,17.0) {$1s$};
    \draw[txh dim] (0,-15.5) -- (16,-15.5); \node[txh lab, anchor=north] at (8,-15.8) {$4a_0$};
  \end{scope}
  \begin{scope}[shift={(60,18)}]
    \foreach \px/\py in {-4.0/4.5,0.9/4.9,3.7/3.6,2.2/-6.2,3.2/1.8,1.2/-7.9,-2.7/3.5,-0.8/-0.4,4.5/1.3,3.4/4.4,10.1/-6.7,-1.2/-0.1,-5.5/-1.7,-3.5/-7.1,-2.3/-2.4,-6.5/6.7,7.5/1.3,-1.0/-5.0,-1.4/-0.6,3.5/2.0,5.3/-1.7,-0.9/-4.9,-1.3/5.1,-5.3/-0.2,-1.0/5.5,7.8/0.6,0.3/1.2,6.0/-4.9,2.2/-4.1,6.8/6.4,-5.2/-3.3,-2.4/3.1,2.9/-10.5,2.1/3.0,-6.9/-8.6,0.8/-6.6,6.9/-1.2,-0.6/-4.3,0.6/0.9,2.1/-9.3,-6.4/-8.2,-3.4/-0.8,-4.1/1.3,-6.1/-3.6,-0.5/-0.2,-2.3/4.3,0.7/-0.5,-2.4/-11.5,-4.2/-1.1,-0.8/-3.5,11.7/-3.1,0.2/8.9,3.7/-9.9,6.8/9.3,-3.8/-1.0,-6.0/6.4,-0.8/5.6,8.1/-1.1,-0.6/-0.7,-3.6/-4.1,3.9/-8.9,-3.9/-0.2,-0.6/0.4,-0.2/7.6,7.8/-3.1,0.7/-1.5,-8.4/-1.0,7.7/-4.5,5.9/2.9,1.1/1.3,2.9/-3.1,1.9/6.0,-2.1/3.6,5.0/-2.6,7.2/3.0,3.7/-4.2,4.9/-2.2,0.0/-0.4,0.3/-1.2,2.4/4.9,7.0/-8.0,6.1/-5.2,3.1/-5.9,-7.4/2.2,-0.3/1.3,-4.8/2.9,-5.8/8.4,10.0/4.7,-10.3/1.8,0.7/-0.9,4.9/-1.3,-0.1/-0.2,4.7/-3.3,-4.5/-10.6,8.6/8.8,1.9/-5.5,1.0/-0.5,0.5/-6.9,-4.5/1.5,0.6/0.0,-0.1/0.2,0.4/-3.9,3.3/6.8,-9.2/-0.2,4.4/-3.5,-5.7/1.1,-5.0/2.7,-7.4/-1.1,-0.1/-0.5,-2.4/-6.0,0.1/0.1,3.5/5.6,0.3/4.1,-4.3/2.1,0.2/0.7,6.6/-0.4,4.7/2.3,2.7/-6.0,-2.8/-5.4,-3.7/-2.6,-5.1/-6.8,-2.5/4.3,-2.5/-3.2,-9.2/7.8,1.1/3.2,0.7/-11.8,-4.7/2.4,-8.5/3.9,-0.9/5.9,-0.5/0.9,-4.6/0.4,0.2/0.6,-0.7/9.4,1.0/-0.0,-5.1/-5.2,5.5/-6.8,0.3/-0.2,-1.2/-3.4,4.3/-0.7,-4.7/-0.5,-0.8/0.5,-3.2/-2.4,-1.5/-7.2,3.6/-5.9,0.1/-6.7,4.3/-1.5,0.2/2.2,0.8/5.0,6.5/-2.0,-8.3/5.0,-5.1/-6.5,1.2/-4.4,7.5/-1.0,-0.1/2.9,-0.8/-10.7,-5.1/7.6,-4.8/-0.3,7.5/-1.9,-0.3/-8.4,-4.1/-6.0,-7.2/-4.0,0.5/-0.4,1.3/5.0,-8.5/-2.5,4.5/-3.1,0.2/5.3,-5.1/-3.6,-2.7/-13.3,-0.6/0.3,-5.2/4.9,0.5/4.0,0.2/1.0,-5.6/-0.3,-1.5/-4.2,3.8/1.3,-0.3/-8.1,2.2/-4.1,5.5/3.5,-8.7/-5.9,-1.0/5.7,6.5/-7.8,1.1/-6.2,5.5/1.0,2.3/2.0,-0.0/0.3,1.7/-5.3,-4.4/2.9,5.2/-4.0,-4.7/-0.6,-2.7/-1.6,-3.3/-5.2,-4.1/-5.7,5.0/-1.3,-6.3/-1.2,2.3/9.3,-4.3/-0.8,-2.3/-6.3,7.9/-1.8,1.1/-3.7,-1.5/-5.9,4.1/1.0,-4.7/3.4,5.6/-1.3,5.8/-1.4,-9.8/0.9,6.4/-6.9,-9.6/-3.0,0.2/-0.2,5.5/2.7,1.0/5.0,5.7/-1.6,4.6/-2.7,0.9/-5.1,2.3/-3.7,-2.3/-6.0,-0.5/0.6,5.0/0.2,3.6/-1.3,2.9/3.3,-0.1/-0.1,1.2/0.5,-7.6/-0.4,0.3/1.4,-1.1/0.0,-10.5/2.3,-0.6/8.1,3.9/-4.0,1.7/-2.8,-6.4/-5.9,-0.5/-0.3,-14.4/-0.8,-1.9/6.5,3.2/6.3,-9.6/-5.5,5.8/1.1,-3.9/3.9,-0.9/-1.6,-7.1/-0.9,-6.2/-1.8,-6.4/1.8,0.0/-0.4,-0.6/-1.0,-4.0/-2.0,0.3/-0.3,9.9/-1.7,-0.2/5.5,4.4/-5.9,-0.4/-0.5,0.4/0.2,6.7/-3.1,5.5/-0.8,-9.8/-1.0,11.9/-0.1,-6.9/1.9,-0.4/0.6,-2.7/4.1,-3.0/-2.5,-1.8/4.2,2.8/-4.3,-0.6/-7.5,8.0/-7.8,-6.9/0.5,-0.9/2.5,-0.5/-3.7,-3.3/-3.4,-0.3/6.9,-2.3/-6.3,3.3/5.0,-7.6/7.7,1.2/0.1,-1.6/-3.7,3.3/4.6,-8.6/2.1,4.0/-7.1,-5.9/1.4,-1.9/-10.1,3.8/1.8,0.7/0.6,-8.0/-5.6,-1.1/0.1,-2.0/7.6,2.0/4.0,2.4/-8.1,-3.0/-4.3,3.4/-2.7,-0.2/1.3,-0.5/-5.5,-1.5/-3.5,-6.7/4.1,1.9/-7.4,-10.4/-0.0,5.2/-0.0,6.4/-3.8,4.9/1.1,2.8/2.1,0.1/0.4,-1.0/0.7,-5.1/-2.7,0.5/-0.3,4.2/-3.3,-0.3/5.9,3.8/-2.2,-1.8/5.7,-3.5/3.5,7.8/-3.3,-1.1/-0.4,6.4/-3.7,-2.2/-2.8,-7.2/-5.8,-0.2/-6.8,3.8/-9.5,5.5/-1.8,3.8/-4.7,-4.3/0.4,4.2/1.1,-6.0/-0.0,3.4/5.1,-2.2/-5.0,-3.6/-5.6,4.5/8.9,-4.3/-5.4,0.6/0.4,4.0/5.0,-0.3/0.1,4.8/1.4,-3.6/-3.2,-5.8/4.6,0.5/-0.6,7.5/3.2,11.3/-3.8,-3.0/-6.2,4.7/6.7,14.6/-0.3,-1.4/4.0,1.0/3.7,-5.8/6.6,-0.5/-0.0,-0.0/-0.2,-3.8/1.8,4.8/2.5,1.3/-5.3,-9.1/-5.0,-6.6/1.2,-0.5/0.8,-10.8/13.0,5.2/-3.5,0.3/0.2,-5.7/-4.2,-2.9/-8.5,0.3/0.2,-9.3/5.4,-4.2/-2.2,0.7/-10.1,-4.8/-4.2,-2.5/-1.1,-11.4/-8.7,3.7/-7.5,-7.3/6.4,-6.5/-0.1,0.5/-8.4,-4.8/6.7,2.8/7.4,1.1/-4.0,-11.5/-2.0,2.7/-5.5,0.3/1.3,-4.5/-0.5,-3.9/12.2,-0.4/5.2,3.2/-4.5,0.3/-9.3,0.2/0.3,-6.2/2.8,-0.6/0.7,-0.9/-7.9,0.8/-4.3,-4.4/-1.5,0.4/0.3,-3.8/1.8,-1.3/-2.5,-6.1/4.1,-1.5/0.2,0.2/-0.5,-3.4/1.9,5.6/-5.5,4.4/7.7,0.9/11.8,3.6/3.1,-2.3/6.2,5.2/8.9,6.8/-2.0,-1.3/-6.7,0.1/-1.6,-9.4/-2.7,-2.4/5.2,7.4/1.4,-0.9/-4.4,5.2/2.9,-5.9/-2.5,-0.7/4.0,-3.1/5.5,4.9/3.0,-0.8/-6.2,-0.8/0.0,10.4/-5.5,-3.0/7.6,1.3/-3.3,0.2/-5.8,-0.1/-1.2,0.8/4.4,-3.4/-11.9,-4.3/-0.4,-0.6/-1.0,8.5/-2.5,-3.4/1.7,1.6/8.8,-1.0/4.2,-5.1/7.2,-2.3/4.8,1.0/-0.9,1.9/-0.1,3.2/-6.9,-9.1/1.0,6.6/1.3,-8.1/-1.1,-0.6/-0.4,-5.6/3.0,6.9/1.9,5.1/3.0,-4.6/3.3,-0.0/-1.0,-5.9/-4.1,-0.5/6.5,-0.1/0.1,5.3/-1.1,0.4/5.6,1.0/1.0,-8.1/0.4,-1.7/-3.6,3.4/3.6,-4.8/2.4,-2.5/-2.5,4.1/-2.2,-2.2/-4.7,4.7/-8.9,0.8/7.6,-3.6/3.1,1.0/-6.6,3.1/-4.2,-3.9/-1.7,-4.9/-3.9,2.1/-9.7,5.2/-8.1,0.1/-0.6,-13.2/-2.0,-6.0/4.1,3.7/10.9,7.7/0.4,-5.7/7.0,-4.5/1.6,1.3/1.4,-2.0/-2.3,5.5/0.9,-4.8/-7.2,5.2/4.2,7.7/8.0,-3.9/-0.5,-0.1/-11.4,-5.4/5.4,0.2/-0.3,7.9/-8.8,1.4/-5.6,-0.1/-3.7,-12.0/0.3,-2.5/3.2,-0.7/0.1,-5.6/5.1,-1.9/2.9,5.0/-5.6,-0.4/-0.0,2.0/7.3,-8.1/1.2,-0.8/-1.3,-7.2/1.3,-8.1/-1.6,-5.1/3.7,6.9/6.3,4.9/-1.6,-0.1/-0.6,-2.9/4.3,-5.2/-2.4,4.6/-5.9,-1.4/4.7,-3.2/-7.5,3.4/-7.6,6.8/2.0,-0.2/-5.9,-3.8/-2.6,-1.7/-7.2,8.3/-0.5,-0.3/-0.3,0.3/-0.6,1.3/-9.3,0.5/-11.7,5.9/3.4,-3.4/-4.3,0.0/1.3,4.3/0.7,-4.3/-4.7,1.9/4.8,4.7/-0.5,0.3/-6.2,-2.6/9.3,0.1/1.2,4.0/-1.4,-4.2/6.4,-5.4/-3.5,-0.9/-0.1,6.2/-1.1,2.9/3.5,3.7/-8.4,-3.9/6.7,-4.0/0.9,0.1/-0.9,2.5/-4.5,2.0/8.9,2.4/11.4,-7.3/-2.1,-8.5/1.1,-1.2/1.1,-0.1/-0.2,6.4/-0.0,-3.7/-5.1,-0.5/-1.3,1.0/-4.3,-5.3/3.8,-6.7/-3.1,1.7/-9.2,-6.8/-1.8,-5.7/-3.5,-2.9/4.9,7.8/-1.0,6.0/6.0,3.2/-3.2,-4.7/-4.8,-4.5/-1.6,-1.2/-6.4,4.5/-1.1,0.6/-6.3,0.0/-0.4,5.7/4.6,0.2/6.9,-0.2/-11.0,-5.0/-2.8,6.0/4.7,-5.8/-0.8,-2.8/-6.1,-6.0/0.0,4.0/-5.5,5.6/4.3,7.6/-2.3,2.6/2.2,5.8/2.4,0.8/-1.3,11.3/4.5,-0.7/0.2,0.1/-0.4,0.7/-0.7,3.3/-8.7,-8.3/3.8,0.4/0.2,0.7/1.2,-1.0/-6.9,-0.1/-6.2,-4.8/-3.1,-3.6/4.4,-4.1/-3.3,-2.2/-0.9,-3.3/4.3,1.1/3.0,3.2/-5.6,6.4/1.1,-0.9/-0.5,6.3/1.2,3.5/-2.9,-2.7/-5.4,-0.6/-0.1,-4.7/4.3,-4.1/-1.2,3.5/1.8,-7.5/-2.4,-2.3/9.1,2.1/-9.2,-1.6/5.1,-3.7/4.7,5.2/-0.1,8.3/1.2,-3.1/3.0,-3.2/10.5,-5.4/0.8,5.7/-1.5,-0.0/0.1,0.2/-0.9,-5.4/-6.2,-0.7/-0.6,-12.6/4.4,-2.9/-2.1,-2.4/-6.0,6.2/-4.0,0.8/-5.7,0.1/4.3,0.8/-0.0,1.0/-6.3,-0.2/-0.7,5.3/2.1,-5.0/1.3,1.8/-7.9,-0.9/-6.5,4.7/2.5,2.8/6.6,8.0/-3.3,0.3/0.4,-4.2/7.9,3.4/5.0,-0.9/-4.0,-2.7/3.7,9.8/-0.7,1.3/-0.9,-5.0/-1.8,1.2/0.6,-0.3/1.6,2.4/-3.0,2.1/-6.0,2.4/-8.3,0.8/6.2,4.5/3.9,1.9/5.3,4.8/3.2,-1.8/-4.9,-1.5/4.4,6.2/8.5,5.8/11.0,-3.0/2.7,-3.4/-9.0,0.0/-1.3,-1.4/7.7,5.6/-4.2,2.4/1.7,0.8/1.0,-0.6/-5.6,-0.1/-0.5,6.3/-0.1,-1.3/3.5,-5.2/-11.0,0.6/4.0,6.2/-0.4,-5.3/-5.3,-4.0/-0.6,-6.5/2.3,4.2/-6.0,-2.2/5.6,1.9/-10.8,1.0/-3.6,3.5/2.3,-4.4/0.1,-3.9/-5.2,-6.6/-5.9,-2.1/-0.2,-2.7/-3.4,-6.3/1.0,5.9/6.5,-6.3/-3.3,0.6/-5.9,-2.5/-2.7,0.1/0.4,7.0/3.5,7.8/11.5,-9.7/-7.9,-0.4/-0.9,8.3/0.9,9.3/-4.9,5.4/-0.1,2.4/-10.2,-2.9/6.0,4.7/2.4,6.8/2.3,-3.3/2.3,11.7/-1.6,-0.0/-6.9,9.0/-1.4,0.5/0.5,4.6/4.6,-6.9/-1.7,-9.4/-3.2,1.1/5.7,-9.4/5.0,-2.1/-9.4,7.7/2.7,5.5/8.0,0.2/-0.9} {\fill[bkBlue!80!black] (\px,\py) circle (0.32);}
    \fill[bkRed!80!black] (0,0) circle (0.8);
    \node[txh lab, anchor=south] at (0,17.0) {$2s$};
    \draw[txh dim] (0,-15.5) -- (16,-15.5); \node[txh lab, anchor=north] at (8,-15.8) {$13a_0$};
  \end{scope}
  \begin{scope}[shift={(102,18)}]
    \foreach \px/\py in {2.0/-6.5,2.2/4.6,6.6/-0.9,-0.7/1.6,0.3/-5.3,0.8/-4.0,8.0/5.6,1.7/-0.8,2.2/5.6,-0.3/6.6,-1.2/3.1,1.2/6.4,3.3/-7.8,0.7/-5.0,1.8/10.1,-6.6/-6.4,2.8/4.4,-1.1/4.4,0.9/-3.2,4.8/5.0,-0.4/4.2,0.0/-1.0,0.2/-0.8,1.4/6.8,0.4/1.8,-4.5/3.2,7.4/5.8,-0.4/-0.6,-0.2/-7.7,-0.0/-7.6,5.9/-12.5,-0.8/1.0,2.7/-4.6,-3.6/3.3,3.2/-4.5,-0.1/-3.4,-0.9/-5.5,-1.3/-3.6,-1.6/-4.6,1.5/-6.0,-2.3/-3.2,1.6/-1.9,-2.2/3.6,0.9/-4.5,4.2/11.9,-7.6/-7.2,2.5/2.7,1.2/7.2,0.2/1.5,2.5/1.0,-2.5/5.6,2.4/-4.6,0.7/-3.8,-3.4/9.4,-2.7/3.7,-2.4/2.3,-0.4/-7.7,2.9/4.7,0.2/-3.8,-0.8/3.4,1.1/-6.2,-2.3/11.6,-1.2/1.4,0.5/-3.2,2.3/-3.5,1.0/3.6,-2.6/2.4,0.9/-5.4,-0.1/3.0,-1.0/2.0,-0.2/6.3,2.1/1.6,-1.5/-4.1,-1.5/1.5,1.7/6.6,-3.4/-4.3,-4.3/3.8,0.5/2.7,-0.8/2.2,-1.3/-3.7,3.8/-2.5,-0.2/-3.8,-0.9/-5.7,-1.2/5.5,1.6/3.2,0.6/-1.5,4.7/-5.3,-0.7/3.4,-5.8/7.9,0.4/-8.1,0.8/-3.1,-2.0/-3.6,1.8/-2.0,0.5/-3.2,-0.0/-1.5,-1.3/-2.5,3.9/1.1,-1.1/-4.0,1.0/-1.9,0.4/4.5,-0.8/-5.2,0.7/5.3,0.4/1.7,-0.8/7.2,-1.1/-5.8,4.3/-6.5,-2.3/-4.4,-0.3/3.0,0.3/-2.6,1.3/-6.3,-0.7/1.8,1.9/5.4,0.5/-3.9,-2.2/-4.2,2.1/-1.9,-4.1/8.1,-0.2/-2.8,-1.4/-2.3,-1.5/5.9,2.0/9.1,1.2/11.4,0.4/4.5,-0.4/-4.4,1.4/-10.8,0.8/-9.0,-0.1/-4.2,1.8/-3.3,-4.0/-10.9,-2.6/1.3,1.2/4.5,-0.4/-3.7,0.2/-5.0,-5.6/6.8,1.9/-5.4,-1.8/1.5,1.1/-2.4,0.6/-2.3,1.5/4.9,1.7/-4.2,2.0/9.7,0.5/-3.2,-2.7/-8.4,-1.7/10.0,-0.6/3.7,-0.0/1.7,0.7/8.0,-3.5/-2.9,1.5/-2.5,3.5/4.4,-0.5/-1.2,1.0/2.5,-3.4/6.2,-3.1/-5.4,1.2/9.2,-0.3/-2.3,-0.7/1.2,-2.0/8.9,-0.2/3.5,-0.3/6.3,-2.7/-5.1,0.1/-6.7,-10.9/10.0,-2.1/-5.0,5.9/4.4,-1.0/4.0,-0.3/-1.9,2.9/-8.1,0.7/-3.9,2.1/-3.6,4.3/-9.8,-1.1/-2.1,-1.0/-6.2,-0.5/-3.0,3.3/3.0,0.5/-3.4,4.3/5.7,2.4/2.3,-0.3/3.2,3.1/0.6,-5.2/2.0,-0.1/3.5,-0.3/4.5,0.1/-3.5,-2.2/-5.6,-2.9/-2.4,5.9/1.5,-1.1/2.1,3.4/-6.3,-2.0/-2.6,1.2/5.2,1.2/-3.8,1.1/-1.8,-1.3/4.1,-0.5/4.4,0.6/7.0,1.8/-3.2,1.3/-2.8,-1.1/-3.6,2.7/5.2,-0.9/-1.7,-2.4/4.5,-0.5/-2.3,-3.2/-4.2,0.1/-7.1,-6.2/3.3,0.8/-4.7,-3.8/-7.1,-2.6/-2.1,-1.6/-3.6,-0.6/-4.2,0.9/-3.1,-0.4/1.3,-1.6/4.0,1.8/4.5,3.7/2.2,4.7/-2.1,-0.7/6.3,1.3/-2.2,-1.6/3.9,1.7/6.4,0.6/-8.9,-1.2/5.8,0.1/-2.1,1.3/-4.0,0.7/-4.7,-0.7/-1.5,2.5/2.6,0.0/-4.5,0.3/-6.0,0.9/-4.0,-3.4/2.8,1.5/2.6,-1.1/-3.2,-0.6/2.7,-1.2/-3.9,-2.9/2.0,-3.3/3.2,-6.4/-2.9,-0.8/-5.0,0.7/-7.2,4.0/2.5,-0.5/-3.5,-3.1/5.0,-0.1/-4.2,1.1/-9.1,-0.6/1.8,-2.7/-9.5,-1.5/3.9,0.4/2.0,2.2/-5.9,-2.4/8.5,0.3/0.8,-0.5/3.3,-2.0/-4.8,-3.3/3.7,-1.2/-1.4,1.5/5.7,0.2/4.8,-4.4/7.4,-0.5/1.8,-1.6/-6.6,2.1/4.9,6.6/2.0,-3.4/2.4,-0.5/5.6,-0.0/-2.7,2.0/1.8,5.3/-7.5,0.7/-4.2,-1.3/0.7,3.7/-6.0,-5.7/-4.5,-1.8/-2.1,-1.5/1.3,1.2/2.3,3.4/4.7,0.1/-1.3,-3.0/-2.0,2.9/7.3,-2.3/-6.6,2.1/-2.5,-1.3/-2.7,-4.1/-4.1,3.2/-1.5,-2.0/-3.3,-0.4/5.0,3.5/6.9,-4.8/-6.8,-0.9/-5.8,1.5/-5.4,-4.3/1.9,0.8/0.8,3.0/2.7,0.4/2.2,-0.8/3.3,-0.4/3.9,0.3/5.8,-1.5/-7.2,0.0/1.4,-0.3/-1.4,1.1/5.9,1.3/-4.1,-3.8/4.1,-2.9/-2.5,-1.8/3.4,0.1/4.1,5.3/-6.1,1.3/-7.8,0.3/4.4,-0.4/4.8,-0.5/3.4,-4.3/4.5,2.1/1.8,-0.8/-2.4,3.1/-11.7,-5.1/-6.4,0.4/3.6,-2.7/3.3,0.5/-1.9,0.7/4.2,7.9/-9.1,-1.0/-2.9,-0.9/-4.1,0.2/6.2,-2.0/-4.2,1.8/-5.4,-0.8/-2.4,-0.6/-1.5,0.2/-3.2,1.0/-2.0,0.4/-4.5,1.1/2.2,-2.4/5.3,3.7/2.5,4.2/-2.5,3.3/4.1,1.5/1.9,-1.4/3.9,-1.5/-3.5,-5.3/-0.4,2.8/-6.8,0.2/-0.8,5.3/-4.1,-3.1/2.3,-2.2/-3.4,-4.8/-3.7,1.0/-4.3,-0.8/-2.1,-3.3/-3.4,2.4/9.6,4.2/6.0,1.1/3.4,0.4/3.6,1.4/-0.8,-3.2/-9.3,5.3/-2.4,-3.2/6.4,0.2/-4.1,0.8/-3.2,5.5/-8.2,-3.4/4.9,-1.7/-2.4,-3.5/5.5,1.4/3.2,2.7/14.0,-5.9/-8.0,0.9/-3.5,0.3/-3.2,-4.3/5.8,5.7/3.2,2.9/-4.8,1.6/6.1,-1.0/-5.0,-0.3/4.2,-0.0/-4.0,-1.5/2.0,-1.7/2.9,0.9/-2.7,1.0/3.3,4.6/-4.3,4.6/-4.4,-0.1/6.7,-3.9/-2.2,1.0/-2.2,-3.1/-2.5,2.7/4.2,1.4/-2.8,0.9/3.2,3.1/-5.4,-1.3/-3.0,-5.6/-5.2,10.5/6.0,0.5/1.7,4.6/6.0,-2.4/3.1,-1.7/-8.6,3.3/-4.8,-0.5/-5.2,2.4/2.8,-1.4/1.6,0.5/-2.6,-0.5/-3.8,1.7/-2.3,4.9/3.3,1.2/-4.6,1.6/-1.9,2.6/-2.9,1.1/-4.5,-3.4/-1.2,-5.3/-7.2,0.2/-5.8,-0.0/-6.6,0.7/-7.6,2.2/-4.8,0.8/3.9,-3.5/-5.2,1.4/2.9,0.7/-7.4,0.4/-4.4,1.6/3.7,3.1/-3.3,0.5/3.9,6.3/2.6,1.2/7.2,3.6/-3.7,-2.8/5.3,-1.5/2.8,2.0/2.8,0.4/-3.5,-1.4/-5.9,-0.1/3.0,1.7/-2.6,0.6/-1.9,-1.1/-2.8,-0.8/2.0,4.7/11.2,-0.6/-3.4,3.2/-2.9,1.7/4.6,3.2/7.0,1.2/3.7,2.0/5.2,3.7/4.2,-1.0/0.9,-1.6/-1.9,-1.6/-4.1,0.5/-1.4,2.0/-5.4,-9.1/-3.0,0.7/2.1,-1.8/5.0,2.4/-2.2,1.8/-5.6,0.3/-4.3,-1.6/6.5,0.2/3.2,-3.6/8.8,1.1/-3.6,-0.6/-2.0,2.0/-0.6,-1.2/-10.6,0.6/-1.3,7.4/-2.3,0.2/-11.8,-3.6/-7.0,-6.9/5.0,-1.1/10.3,-0.2/-4.3,2.7/7.9,-3.3/-2.5,6.7/-8.4,2.0/3.9,-0.7/0.4,6.5/-5.6,-4.2/-5.2,3.2/4.4,2.2/5.1,0.6/5.5,0.2/-3.5,-1.3/-6.5,-0.1/3.2,-2.3/3.4,4.0/3.0,0.9/3.0,-2.6/4.3,-0.3/8.2,-3.3/2.3,1.2/2.9,1.0/8.2,0.6/5.1,1.5/-1.5,0.1/3.2,1.0/-2.9,-1.7/4.4,-3.0/-3.6,-1.8/-4.8,1.9/-2.0,1.5/-3.9,-6.8/5.9,-2.8/-7.4,-5.6/7.9,1.7/1.5,5.1/2.4,-1.7/-4.5,2.3/4.5,-6.2/-8.2,0.1/9.8,-0.6/-1.5,-2.8/-7.0,3.6/4.3,3.1/2.3,2.4/3.1,0.9/3.5,0.9/-0.9,-2.5/3.7,1.2/3.1,0.8/7.4,-0.5/3.3,0.8/2.9,-2.3/2.5,-0.6/3.7,2.6/1.8,-2.1/-4.4,-1.0/6.7,-2.3/1.8,-2.3/-4.1,-0.3/4.4,-0.2/-1.6,-2.8/2.5,2.2/-2.7,-8.3/7.7,-2.7/2.7,-1.8/4.2,3.9/-2.0,-1.1/9.1,0.5/-7.4,3.0/3.2,-1.2/2.9,-1.6/-3.0,0.8/-2.0,-2.2/-1.6,-3.3/5.4,-1.8/-3.0,-2.3/1.1,1.4/3.3,-0.6/2.4,1.5/-4.6,1.2/2.2,-0.7/1.4,3.5/3.6,3.1/3.4,-1.2/1.8,4.6/-2.6,-0.8/-3.8,3.0/6.7,1.2/3.7,1.7/-1.5,4.7/7.4,0.6/3.1,-0.7/6.1,3.7/10.3,-1.6/-4.4,-0.6/4.7,0.3/-5.6,-0.5/8.9,-0.3/-6.2,-5.7/-5.3,0.1/-1.3,1.2/3.6,0.4/1.7,0.6/2.9,-0.1/-5.0,0.2/3.5,-1.9/-3.7,-0.5/-1.9,0.3/-2.9,1.9/5.5,-1.5/3.4,0.2/2.1,1.5/2.8,-4.6/9.1,1.1/-3.6,2.5/-1.5,1.7/0.3,-4.3/6.0,2.6/4.9,-1.1/2.2,-1.3/3.2,-1.4/5.5,-2.5/2.1,-4.1/-3.5,4.1/7.2,-0.0/6.2,-3.7/-5.8,-0.7/-1.8,-1.1/-5.6,-1.3/-1.8,-2.6/11.3,-0.5/-3.3,-0.9/-3.7,-1.9/3.7,2.1/-5.2,-1.5/-2.9,1.2/-3.8,5.4/2.2,0.5/-10.6,0.7/7.0,-0.4/-4.3,1.8/-6.0,1.0/-4.9,3.4/2.9,-4.3/-7.5,-1.7/3.8,-2.5/4.4,1.7/-3.7,-3.7/1.2,-2.2/2.4,0.5/-6.3,0.8/1.9,-1.2/-1.7,-2.4/-5.3,-5.1/4.3,-2.8/3.2,-2.6/4.0,2.2/2.9,-0.3/2.3,-2.1/6.4,-0.1/-1.7,-1.7/-3.0,-1.3/-2.9,-2.9/-2.3,1.0/3.6,0.4/6.3,-0.3/3.7,-4.7/7.7,0.8/-8.4,-2.2/-8.3,2.7/4.2,1.6/-11.4,0.7/-2.8,-2.1/-6.2,-0.7/-3.3,-1.2/2.9,0.4/2.3,-0.7/4.5,5.5/-2.2,-2.5/-5.7,0.7/-6.3,-2.4/-7.5,0.9/-2.5,-0.2/-4.2,2.3/-1.9,-6.1/-9.9,-0.6/-3.2,-3.4/2.6,2.6/8.3,2.6/5.1,2.2/4.6,2.7/4.3,3.3/-3.5,-1.5/1.7,2.7/-7.4,0.4/4.7,-7.5/-3.1,2.4/-2.4,1.0/-1.4,-1.2/2.6,0.0/-3.4,-4.6/4.7,-2.7/-4.2,2.9/-0.9,-0.9/-3.6,1.4/7.9,-5.3/8.4,1.0/-2.9,0.0/-1.7,0.9/-3.9,-1.5/3.8,-0.5/4.2,0.7/-4.8,-1.5/-1.3,4.2/-1.8,0.3/2.7,-0.0/-1.7,-4.2/2.2,-1.9/4.8,0.0/1.7,0.6/-1.6,2.7/-3.7,-0.8/1.7,0.2/-5.4,7.7/-3.4,1.4/4.1,-0.9/-4.9,2.0/5.4,1.5/-4.0,2.6/1.7,-4.6/-8.0,-4.0/-6.6,0.8/5.6} {\fill[bkBlue!80!black] (\px,\py) circle (0.32);}
    \fill[bkRed!80!black] (0,0) circle (0.8);
    \node[txh lab, anchor=south] at (0,17.0) {$2p$（たての向き）};
    \draw[txh dim] (0,-15.5) -- (16,-15.5); \node[txh lab, anchor=north] at (8,-15.8) {$13a_0$};
  \end{scope}
\end{tikzpicture}
\bkfigcaption{水素原子の電子の雲（原子核を通るうすい板の中だけを、点で打った。点1つが、1回の測定で電子が見つかった場所。まん中の丸が原子核）。$1s$ と、$2s$・$2p$ では、目もり（下の矢印の長さ）が違う。$2s$ の雲の中には、点のない輪（節）がある。$2p$ は、まん中を通る横の面の上には電子が見つからない。}\label{ph:原-7:fig-cloud}
\end{bkfigure}

雲の形には、名前がついています。丸い雲が $s$、ひょうたんの形が $p$、さらに複雑な形が $d$・$f$ です（もとは、光のスペクトルの線の見え方の英語の頭文字です）。数学の本（4-15）で約束したように、雲の形は、原子核を原点にした球座標で書くと、とても簡単な式になります。$1s$ は、原子核からの距離 $r$ だけで決まり、$\theta$ や $\varphi$ によりません。だから丸いのです。$2p$ は $\cos\theta$ に比例する部分をもつので、たての向き（$\theta=0$ と $180^\circ$）に長く、横（$\theta=90^\circ$）では0になります（⑩）。

\begin{phmijika}[雲の「大きさ」はどれくらい？]
$1s$ の雲は、原子核から約 \qty{0.16}{nm}（ボーア半径 $a_0$ の3倍）までに、電子が見つかる確率の94\,\%が入ります（例題2 と同じ計算）。それより外にも、確率は小さいながら0ではなく、雲にははっきりしたふちがありません。原子の「大きさ」は、こうしたぼやけた雲の大きさなので、はかり方によって少し違う値になります。
\end{phmijika}

\phstage{te}{電子の席を数える}{3}

\Needspace{12\baselineskip}\bksubsection{3}{水素原子のエネルギーと、1s の雲}
\label{ph:原-7:sec3}
水素原子は、陽子1個の周りに電子が1個だけある、いちばん簡単な原子です。電子は、陽子からの距離 $r$ で、クーロンの力の位置エネルギー
\pheq{U(r)=-\frac{ke^2}{r}}{J ＝ (N·m$^2$/C$^2$) × C$^2$ ÷ m}
をもちます（電-1・電-2。$k=8.99\times10^9$ N·m$^2$/C$^2$）。これは、陽子に近いほど深い「すりばち」の形の井戸です。原-5 の箱や原-6 の井戸と同じように、この井戸に閉じこめられた電子の波は、とびとびのエネルギーの定在波になります。3次元のシュレディンガー方程式を解くと（⑩）、エネルギーは
\pheq{E_n=-\frac{ke^2}{2a_0}\cdot\frac{1}{n^2}=-\frac{\qty{13.6}{eV}}{n^2}\qquad(n=1,2,3,\ldots)}{eV ＝ eV}
\label{ph:原-7:En}となります。$a_0=\hbar^2/(m_{\mathrm e}ke^2)\approx\qty{52.9}{pm}$ はボーア半径です。これは、原-4 ⑤でボーアの模型から求めた値と、\textbf{まったく同じ}です。ボーアの模型は「電子は円の上を回り、1周の長さが波長の整数倍」という仮定から出しましたが、シュレディンガー方程式は、円の軌道を考えずに同じ答えを出します。そのかわり、電子の居場所は、決まった円ではなく雲になります。

\begin{phyomi}[$E_n=-13.6\ \mathrm{eV}/n^2$]
「水素原子の電子のエネルギーは、$n$ の2乗に反比例する負の値」と読みます。負なのは、電子を陽子からうんと遠く（$r\to\infty$、$U=0$）へ引きはなすのに、エネルギーが要るからです。$n=1$ がいちばん低く（\textbf{基底状態}）、電子をそこから取り去るのに \qty{13.6}{eV} が要ります。これを\textbf{電離エネルギー}（イオン化エネルギー）といいます。\bkidx{でんりえねるぎー}{電離エネルギー}$n=2$ は \qty{-3.40}{eV}、$n=3$ は \qty{-1.51}{eV} で、$n$ が大きいほど段の間かくがつまっていきます。
\end{phyomi}

\phboxhead{1s の雲：いちばん見つかりやすい距離}
いちばん低い $1s$ の波動関数は、
\pheq{\psi_{1s}(r)=\frac{1}{\sqrt{\pi a_0^{\,3}}}\,e^{-r/a_0}}{m$^{-3/2}$ ＝ m$^{-3/2}$}
\label{ph:原-7:psi1s}です（⑩で、シュレディンガー方程式をみたすことを確かめます）。$|\psi|^2$ は、その場所の小さな体積 \qty{1}{m^3} あたりに電子が見つかる確率（\textbf{確率密度}）です。前の係数は、全部の空間で確率をたすと1になるように選んであります（原-5 ④の規格化）。$|\psi|^2=e^{-2r/a_0}/(\pi a_0^{\,3})$ は、原子核の所（$r=0$）でいちばん大きく、外へ行くほど小さくなります。

ところが、「原子核からの距離が $r$ と $r+\Delta r$ のあいだに電子が見つかる確率」を考えると、話が変わります。この場所は、半径 $r$、厚さ $\Delta r$ のうすい球の殻で、その体積は $4\pi r^2\Delta r$ です（4-5 の球の表面積×厚さ）。確率は「確率密度×体積」なので、
\pheq{P(r)\,\Delta r=|\psi|^2\times4\pi r^2\Delta r=\frac{4r^2}{a_0^{\,3}}\,e^{-2r/a_0}\,\Delta r}{1 ＝ m$^{-3}$ × m$^3$}
となります。$P(r)$ を\textbf{動径分布}といいます。\bkidx{どうけいぶんぷ}{動径分布}$r$ が小さいと殻の体積 $4\pi r^2\Delta r$ が小さく、$r$ が大きいと $e^{-2r/a_0}$ が小さいので、そのあいだに山があります。$P(r)$ のグラフ（図\ref{ph:原-7:fig-radial}）を描くと、山はちょうど $r=a_0$ の所にあります（微分で確かめる方法は、下のもっと知りたい）。\textbf{電子がいちばん見つかりやすい距離は、ちょうどボーア半径 $a_0$}。ボーアの円の半径は、雲の中で電子がいちばん多く見つかる距離として、量子力学の中に生き残っているのです（原-5 のはてな？の続き）。

\begin{bkfigure}
\begin{tikzpicture}
\begin{axis}[bkplot, axis lines=left, width=112mm, height=50mm, xmin=0, xmax=6.3, ymin=0, ymax=1.15, xtick={0,1,2,3,4,5,6}, xticklabels={0,$a_0$,$2a_0$,$3a_0$,$4a_0$,$5a_0$,$6a_0$}, ytick={0,0.5,1}, xlabel={原子核からの距離 $r$}, ylabel={確率の大きさ}]
  \addplot[bkcurve, domain=0:6.2, samples=150] {4*x^2*exp(-2*x)};
  \addplot[bkcurve2, dashed, domain=0:6.2, samples=150] {exp(-2*x)};
  \draw[bkGray, line width=0.5pt, densely dotted] (axis cs:1,0) -- (axis cs:1,0.5413);
  \draw[bkGray, line width=0.5pt, densely dotted] (axis cs:1.5,0) -- (axis cs:1.5,0.448);
  \node[font=\sffamily\footnotesize, anchor=west] at (axis cs:1.6,0.74) {動径分布 $P(r)$（実線）：山は $r=a_0$};
  \node[font=\sffamily\footnotesize, anchor=west] at (axis cs:0.25,1.06) {$|\psi|^2$（破線。$r=0$ の値を1とした）};
  \node[font=\sffamily\footnotesize, anchor=south west] at (axis cs:1.55,0.46) {平均 $1.5a_0$};
\end{axis}
\end{tikzpicture}
\bkfigcaption{水素原子の $1s$ の、確率密度 $|\psi|^2$（破線）と、動径分布 $P(r)$（実線。$a_0$ を単位にした値）。$|\psi|^2$ は原子核の所がいちばん大きいが、原子核のすぐ近くは殻の体積が小さいので、距離で数えるといちばん多いのは $r=a_0$。距離の平均は $1.5a_0$（例題3）。}\label{ph:原-7:fig-radial}
\end{bkfigure}

\begin{motto}[4]
\textbf{山の位置と、内側にいる確率を、微分と積分で求める}\par
$P(r)$ の山の位置は、$P(r)$ を $r$ で微分して0とおくと求まります（積の微分と $e^{cr}$ の微分：5-6・5-14）。前の定数 $4/a_0^{\,3}$ は山の位置に関係しないので、
\[ \frac{\dd}{\dd r}\bigl(r^2e^{-2r/a_0}\bigr)=\Bigl(2r-\frac{2r^2}{a_0}\Bigr)e^{-2r/a_0}=0\qquad\text{より}\qquad r=a_0 \]
です。また、$P(r)\,\Delta r$ を $r=0$ からある距離まで合わせる（積分する：5-10）と、その距離より内側に電子が見つかる確率になります。$a_0$ より内側は約 0.32（例題2）です。「距離×その距離にいる確率」を合わせると距離の平均になり、$1.5a_0$ です（例題3）。
\end{motto}

\begin{phbaisuru}{原子核の電荷 $Ze$}
\noindent\begin{minipage}[c]{0.45\linewidth}\raggedright
電子1個の原子で、原子核の電荷を2倍（$Z=2$：電子を1個失ったヘリウム He$^+$）にすると、電子は2倍の強さで引かれ、雲の大きさは半分（$a_0/Z$）に縮みます。エネルギーは $-13.6\,Z^2/n^2$ eV で、4倍深くなります（\qty{-54.4}{eV}）。小さく縮んだ分だけ、位置エネルギーがもっと負になるからです。原子番号の大きい原子の内側の電子が、強く引かれて原子核のすぐ近くにいるのは、このためです。
\end{minipage}\hfill
\begin{minipage}[c]{0.52\linewidth}\centering
\begin{tikzpicture}
\begin{axis}[bkplot, axis lines=left, width=50mm, height=40mm, xmin=0, xmax=4.2, ymin=0, ymax=1.2, xtick={0,0.5,1,2,3,4}, xticklabels={0,$\tfrac12$,1,2,3,4}, ytick=\empty, xlabel={$r$（$a_0$ を単位に）}, ylabel={$P(r)$}]
  \addplot[bkcurve, domain=0:4.1, samples=120] {4*x^2*exp(-2*x)};
  \addplot[bkcurve2, dashed, domain=0:4.1, samples=120] {4*8*x^2*exp(-4*x)};
  \node[font=\sffamily\footnotesize, anchor=west] at (axis cs:1.3,0.62) {$Z=1$（実線）};
  \node[font=\sffamily\footnotesize, anchor=west] at (axis cs:0.75,1.08) {$Z=2$（破線）};
\end{axis}
\end{tikzpicture}
\end{minipage}
\end{phbaisuru}

\Needspace{8\baselineskip}\bksubsection{4}{量子数で席を数える：$2n^2$}
\label{ph:原-7:sec4}
1次元の箱の電子の定在波は、腹の数 $n$ という1つの数で区別できました（原-5 ③）。3次元の原子の中の定在波は、3つの向きにゆれ方があるので、区別するのに3つの数が要ります。さらに、電子がもつ「スピン」の向き（⑥）で、もう1つ数が要ります。この4つの数を\textbf{量子数}といいます。\bkidx{りょうしすう}{量子数}

\begin{bkcenter}
\begin{bktable}{colspec={Q[l,m,wd=30mm] Q[c,m,wd=12mm] Q[l,m,wd=34mm] Q[l,m,wd=40mm]}, rowsep=1pt, row{1}={bg=bkTheme!10}, cells={font=\small}}
名前 & 記号 & とれる値 & 何を決めるか \\
主量子数 & $n$ & $1,2,3,\ldots$ & エネルギーの大まかな段と、雲の大きさ（殻） \\
方位量子数 & $l$ & $0,1,\ldots,n-1$ & 雲の形（$l=0,1,2,3$ が $s,p,d,f$） \\
磁気量子数 & $m_l$ & $-l,\ldots,-1,0,1,\ldots,l$ & 雲の向き（$2l+1$ 通り） \\
スピン量子数 & $m_s$ & $+\tfrac12$ と $-\tfrac12$ & スピンの向き（上向き・下向き） \\
\end{bktable}
\bktabcaption{原子の中の電子の状態を区別する、4つの量子数}\label{ph:原-7:tab-qn}
\end{bkcenter}

\noindent 表\ref{ph:原-7:tab-qn}の決まりの出どころ（なぜ $l$ が $n-1$ までなのか、なぜ $m_l$ が $2l+1$ 通りなのか）は、シュレディンガー方程式を球座標で解くとわかります（⑩で結果を読みます）。ここでは、決まりを使って席を数えます。

量子数の組 $(n,l,m_l,m_s)$ が1つ決まると、電子の状態が1つ決まります。これを1つの\textbf{席}とよぶことにします。⑦で学ぶように、\textbf{1つの席には電子が1個までしか入れません}。だから、ある $n$ の殻に入れる電子の数は、その $n$ の席の数です。数えてみましょう。

\begin{enumerate}[label={\bkcircled{\arabic*}}, leftmargin=2.4\zw, itemsep=1pt, topsep=2pt]
\item 1つの $l$ について、$m_l$ は $-l$ から $l$ まで $2l+1$ 通り。それぞれに $m_s$ が2通り。だから、1つの $l$ の席は $2(2l+1)$ 個です。$l=0,1,2,3$（$s,p,d,f$）で、2・6・10・14 個。これが、周期表のブロックのはばです（①）。
\item 1つの $n$ では、$l$ が $0$ から $n-1$ まで。だから、$n$ の殻の席の数は
\[ \sum_{l=0}^{n-1}2(2l+1)=2\bigl(1+3+5+\cdots+(2n-1)\bigr) \]
です（$\sum$ は 5-10 の「たし算をまとめて書く」記号）。
\item 1から始まる奇数を $n$ 個たすと、$n^2$ になります。図\ref{ph:原-7:fig-odd}のように、1・3・5・…の点を L 字の形に並べていくと、正方形になるからです。
\end{enumerate}
\pheq{\text{$n$ の殻の席の数}=2n^2}{個 ＝ 個}
\label{ph:原-7:2n2}$n=1,2,3,4$（K 殻・L 殻・M 殻・N 殻）で、2・8・18・32 です。化学で習う「電子殻に入る電子の数」は、量子数を数えた結果だったのです。


\begin{bkfigure}
\begin{tikzpicture}[x=1mm, y=1mm]
  \foreach \i in {1,...,4} {\foreach \j in {1,...,4} {
    \pgfmathtruncatemacro\mx{max(\i,\j)}
    \ifcase\mx\or
      \fill[bkBlue!80!black] (\i*6,\j*6) circle (1.5);\or
      \draw[bkOrange!85!black, line width=0.9pt] (\i*6,\j*6) circle (1.3);\or
      \fill[bkGreen!55!black] (\i*6-1.3,\j*6-1.3) rectangle (\i*6+1.3,\j*6+1.3);\or
      \fill[bkPurple!80!black] (\i*6-1.5,\j*6-1.2) -- (\i*6+1.5,\j*6-1.2) -- (\i*6,\j*6+1.5) -- cycle;\fi}}
  \foreach \k in {1,2,3} {\draw[bkGray, line width=0.5pt, densely dashed] ({\k*6+3},3) -- ({\k*6+3},{\k*6+3}) -- (3,{\k*6+3});}
  \draw[bkGray, line width=0.4pt] (3,3) rectangle (27,27);
  \fill[bkBlue!80!black] (34,24) circle (1.5); \node[txh lab, anchor=west] at (37,24) {1（$s$）};
  \draw[bkOrange!85!black, line width=0.9pt] (34,18) circle (1.3); \node[txh lab, anchor=west] at (37,18) {3（$p$）};
  \fill[bkGreen!55!black] (32.7,10.7) rectangle (35.3,13.3); \node[txh lab, anchor=west] at (37,12) {5（$d$）};
  \fill[bkPurple!80!black] (32.5,4.8) -- (35.5,4.8) -- (34,7.5) -- cycle; \node[txh lab, anchor=west] at (37,6) {7（$f$）};
  \node[txh lab, anchor=north west] at (3,1) {$1+3+5+7=16=4^2$};
  % 席の箱（n=2）
  \begin{scope}[shift={(80,6)}]
    \node[txh lab, anchor=south west] at (0,15) {$n=2$ の席（箱1つに $m_s$ が2つ）};
    \txhbox{(0,4)}{0}{0}
    \txhbox{(14,4)}{0}{0}\txhbox{(20,4)}{0}{0}\txhbox{(26,4)}{0}{0}
    \node[txh lab, anchor=north] at (3,3) {$2s$};
    \node[txh lab, anchor=north] at (23,3) {$2p$（$m_l=-1,0,1$）};
    \node[txh lab, anchor=west] at (33,7) {$=2\times4=8$};
  \end{scope}
\end{tikzpicture}
\bkfigcaption{左：奇数を順にたすと、正方形の点の数（$n^2$）になる。左下の1点（●）が $s$、次の L 字の3点（○）が $p$、その次の5点（■）が $d$、7点（▲）が $f$ の雲の形と向きの数（L 字の境目は破線）。右：$n=2$ の殻の席。箱1つが $(n,l,m_l)$ の1つの組で、箱には上向き・下向きのスピンが1個ずつ、合わせて2個まで入る。}\label{ph:原-7:fig-odd}
\end{bkfigure}

\begin{phyomi}[$2n^2$]
「主量子数 $n$ の殻には、$2n^2$ 個の席がある」と読みます。2 はスピンの向きの数、$n^2$ は雲の形と向きの数です。この式は、殻の\textbf{席の数}であって、周期表の\textbf{行の長さ}とは違います。第3周期の行の長さは18ではなく8です。なぜかは⑤で見ます。
\end{phyomi}

\phboxhead{表計算で、席を数える}
A 列に $n$、B 列に $l$ を書き、C 列に $2l+1$（$m_l$ の数）、D 列に $2(2l+1)$（席の数）を計算させます。E 列は、同じ $n$ の席の合計です。
\begin{bkcenter}\small
\texttt{C2 = 2*B2+1}\quad\texttt{D2 = 2*C2}\quad\texttt{E2 = SUMIF(A\$2:A\$11, A2, D\$2:D\$11)}
\end{bkcenter}
\begin{bkcenter}
\bksheet{5}{19}{%
1 & $n$ & $l$ & $2l+1$ & 席 & $n$ の合計 \\
2 & 1 & 0（$s$） & 1 & 2 & 2 \\
3 & 2 & 0（$s$） & 1 & 2 & 8 \\
4 & 2 & 1（$p$） & 3 & 6 & 8 \\
5 & 3 & 0（$s$） & 1 & 2 & 18 \\
6 & 3 & 1（$p$） & 3 & 6 & 18 \\
7 & 3 & 2（$d$） & 5 & 10 & 18 \\
8 & 4 & 0（$s$） & 1 & 2 & 32 \\
9 & 4 & 1（$p$） & 3 & 6 & 32 \\
10 & 4 & 2（$d$） & 5 & 10 & 32 \\
11 & 4 & 3（$f$） & 7 & 14 & 32 \\}
\bktabcaption{席の数の表計算。E 列（$n$ の合計）は、どの $n$ でも $2n^2$ になっている。}\label{ph:原-7:tab-sheet}
\end{bkcenter}

\Needspace{12\baselineskip}\bksubsection{5}{電子配置と周期表}
\label{ph:原-7:sec5}
原子番号 $Z$ の原子には、電子が $Z$ 個あります。電子は、エネルギーの低い席から順に入っていきます。どの席に何個入っているかを\textbf{電子配置}といい、$1s^2\,2s^2\,2p^6\,3s^1$（ナトリウム）のように、「$n$ と $l$ の記号」の右肩に、そこに入っている電子の数を書きます。\bkidx{でんしはいち}{電子配置}

\phboxhead{電子が2個以上あると、$l$ でもエネルギーが違う}
水素原子では、エネルギーは $n$ だけで決まりました（$2s$ と $2p$ は同じ高さ）。電子が2個以上ある原子では、そうではありません。外側の電子は、原子核の $+Ze$ の電荷を、内側の電子の雲が一部打ち消した、弱い電荷で引かれます。これを\textbf{遮へい}といいます。\bkidx{しゃへい}{遮へい（原子核の電荷）}$s$ の雲は原子核の近くまでしみこむので（図\ref{ph:原-7:fig-cloud}の $2s$ の内側の小さな雲）、遮へいされにくく、同じ $n$ の $p$ より低いエネルギーになります。$d$ はもっと外側にいて、遮へいされやすいので、$3d$ は $4s$ より少し高くなってしまいます。こうして、席がうまる順番は、ふつう次のようになります。
\[ 1s\to2s\to2p\to3s\to3p\to4s\to3d\to4p\to5s\to4d \]
\[ \to5p\to6s\to4f\to5d\to6p\to7s\to5f\to6d\to7p \]

\begin{bkfigure}
\begin{tikzpicture}[x=1mm, y=1mm]
  \foreach \row/\items in {1/{0/1s/2/bkBlue!14},
                          2/{0/2s/2/bkBlue!14,2/2p/6/bkOrange!16},
                          3/{0/3s/2/bkBlue!14,2/3p/6/bkOrange!16},
                          4/{0/4s/2/bkBlue!14,2/3d/10/bkGreen!14,12/4p/6/bkOrange!16},
                          5/{0/5s/2/bkBlue!14,2/4d/10/bkGreen!14,12/5p/6/bkOrange!16},
                          6/{0/6s/2/bkBlue!14,2/4f/14/bkPurple!12,16/5d/10/bkGreen!14,26/6p/6/bkOrange!16},
                          7/{0/7s/2/bkBlue!14,2/5f/14/bkPurple!12,16/6d/10/bkGreen!14,26/7p/6/bkOrange!16}} {
    \foreach \x/\nm/\w/\col in \items {
      \filldraw[fill=\col, draw=bkInk!70, line width=0.4pt] ({\x*3.4},{-(\row-1)*6.2}) rectangle ++({\w*3.4},-5.2);
      \node[font=\sffamily\footnotesize] at ({(\x+\w/2)*3.4},{-(\row-1)*6.2-2.6}) {$\nm$};}
    \node[txh lab, anchor=east] at (-2,{-(\row-1)*6.2-2.6}) {第\row 周期};}
  \foreach \row/\L in {1/2,2/8,3/8,4/18,5/18,6/32,7/32} {\node[txh lab, anchor=west] at (111,{-(\row-1)*6.2-2.6}) {計 \L};}
\end{tikzpicture}
\bkfigcaption{席がうまる順番と、周期の長さ。1つの行が1つの周期で、箱のはばが、その席の数（$s$ は2、$p$ は6、$d$ は10、$f$ は14）。行の席をたすと、周期の長さ 2・8・8・18・18・32・32 になる。}\label{ph:原-7:fig-order}
\end{bkfigure}

図\ref{ph:原-7:fig-order}のように、$p$ の席（6個）がうまったところで、次の $s$ の席までのエネルギーの差が大きくなります。この大きな段のすぐ下で終わる行が、周期表の1つの周期です。第3周期は $3s$ と $3p$ の8個で終わり、$3d$ の10個は、$4s$ のあとの第4周期に回ります。これが、$2n^2$ では18なのに、第3周期の行の長さが8である理由です。

\phboxhead{3つの原子の電子配置}
\begin{bkcenter}
\begin{bktable}{colspec={Q[l,m,wd=26mm] Q[l,m,wd=50mm] Q[l,m,wd=42mm]}, rowsep=1pt, row{1}={bg=bkTheme!10}, cells={font=\small}}
原子（$Z$） & 電子配置 & 短く書くと \\
ナトリウム（11） & $1s^2\,2s^2\,2p^6\,3s^1$ & $[\mathrm{Ne}]\,3s^1$ \\
鉄（26） & $1s^2\,2s^2\,2p^6\,3s^2\,3p^6\,3d^6\,4s^2$ & $[\mathrm{Ar}]\,3d^6\,4s^2$ \\
銅（29） & $1s^2\,2s^2\,2p^6\,3s^2\,3p^6\,3d^{10}\,4s^1$ & $[\mathrm{Ar}]\,3d^{10}\,4s^1$ \\
\end{bktable}
\bktabcaption{電子配置の例。$[\mathrm{Ne}]$ は「ネオンと同じ内側の電子（10個）」、$[\mathrm{Ar}]$ は「アルゴンと同じ内側の電子（18個）」の意味。}\label{ph:原-7:tab-config}
\end{bkcenter}

\noindent 銅は、順番どおりなら $3d^9\,4s^2$ のはずですが、実際には $3d$ がうまって $4s$ が1個になります。$3d$ と $4s$ のエネルギーはとても近いので、$3d$ をちょうど満席にしたほうが、原子全体として低いエネルギーになるのです。このような例外は、$d$ や $f$ の席のある原子でいくつか見られます。

\Needspace{12\baselineskip}\bksubsection{6}{スピン：向きが2つしかない小さな磁石}
\label{ph:原-7:sec6}
シュテルンとゲルラッハは、銀の原子を細いビームにして、場所によって強さの違う磁場の中を通しました。原子が小さな磁石で、その向きが何でもよいなら、ビームは上下に連続して広がるはずです。ところが、ビームは\textbf{2つにだけ分かれました}（シュテルンは、この分子線の方法などで、1943年のノーベル物理学賞を受けました）。銀の原子の磁石の性質は、いちばん外の $5s$ の電子1個のものです。電子は、それ自身が小さな磁石で、その向きは、磁場に対して「上向き」か「下向き」の2つしかとれないのです。この電子の性質を\textbf{スピン}といいます。\bkidx{すぴん}{スピン}④の4つめの量子数 $m_s=\pm\tfrac12$ は、この2つの向きを表します。

「スピン」は「回転」という意味の英語ですが、電子が本当にこまのように回っているわけではありません（電子の大きさは、はかれる限り0です。回っていると考えると、表面が光の速さをこえてしまいます）。電子が生まれつきもっている、角運動量（回る勢い）と同じ性質の量で、量子力学と相対論を合わせた理論（ディラック）から自然に出てきます。

\phboxhead{磁場の中で、エネルギーが2つに分かれる}
電子の磁石の強さ（磁気モーメント）の目もりは、
\pheq{\mu_{\mathrm B}=\frac{e\hbar}{2m_{\mathrm e}}=\qty{9.27e-24}{J/T}=\qty{5.79e-5}{eV/T}}{J/T ＝ C × J·s ÷ kg}
で、\textbf{ボーア磁子}といいます。\bkidx{ぼーあじし}{ボーア磁子}磁場 $B$ の中で、電子の磁石が磁場の向きにそろうと、エネルギーは $\tfrac12g\mu_{\mathrm B}B$ だけ下がり、反対向きだと同じだけ上がります（磁場の中の磁石は、そろう向きに回る：電-6 ⑪）。$g$ は $g$ 因子とよばれる数で、電子では 2.0023 です（ディラックの理論では2で、わずかなずれは「場の量子論」で説明され、とても精密に確かめられています）。上と下の差は
\pheq{\Delta E=g\mu_{\mathrm B}B}{eV ＝ (eV/T) × T}
\label{ph:原-7:zeeman}です。

\noindent\begin{minipage}[c]{0.48\linewidth}\raggedright
この差のエネルギー $\Delta E$ をちょうどもった電波（光子のエネルギー $hf=\Delta E$：原-2）を当てると、エネルギーの低い向きの電子が電波を吸って、高い向きへ移ります。これを\textbf{電子スピン共鳴}（ESR）といいます。\bkidx{でんしすぴんきょうめい}{電子スピン共鳴}周波数は $f=\Delta E/h$ で、\qty{1}{T}あたり約 \qty{28.0}{GHz}です（例題6）。共鳴する電波の周波数から、物質の中に、相手のいない電子（不対電子）がどれだけ、どんな環境にあるかがわかります。
\end{minipage}\hfill
\begin{minipage}[c]{0.49\linewidth}\centering
\begin{tikzpicture}
\begin{axis}[bkplot, axis lines=left, width=50mm, height=44mm, xmin=0, xmax=2.2, ymin=-1.9, ymax=1.9, xtick={0,1,2}, ytick={0}, xlabel={$B$（T）}, ylabel={エネルギー}]
  \addplot[bkcurve, domain=0:2.1] {0.6*x};
  \addplot[bkcurve2, dashed, domain=0:2.1] {-0.6*x};
  \draw[bkInk, line width=0.8pt, {Stealth[length=1.6mm]}-{Stealth[length=1.6mm]}] (axis cs:1,-0.6) -- (axis cs:1,0.6);
  \draw[bkInk, line width=0.8pt, {Stealth[length=1.6mm]}-{Stealth[length=1.6mm]}] (axis cs:2,-1.2) -- (axis cs:2,1.2);
  \node[font=\sffamily\footnotesize, anchor=east] at (axis cs:0.95,0.05) {$\Delta E$};
  \node[font=\sffamily\footnotesize, anchor=east] at (axis cs:1.95,0.25) {$2\Delta E$};
  \node[font=\sffamily\footnotesize, anchor=west] at (axis cs:0.05,1.68) {上：反対向き（実線）};
  \node[font=\sffamily\footnotesize, anchor=west] at (axis cs:0.05,-1.45) {下：そろう（破線）};
\end{axis}
\end{tikzpicture}
\end{minipage}

\begin{phbaisuru}{磁場 $B$}
上の図のように、上と下の差 $\Delta E$ は $B$ に比例するので、2倍になります。共鳴する周波数も2倍です。MRI の 1.5 T と 3 T では、陽子の共鳴の周波数が 63.9 MHz と 127.7 MHz で、ちょうど2倍です（例題7）。
\end{phbaisuru}

\phboxhead{陽子もスピンをもつ：NMR と MRI}
陽子（水素の原子核）もスピンをもち、磁場の中で向きが2つに分かれます。ただし、陽子の磁石は電子の約660分の1の強さしかなく、共鳴の周波数は \qty{1}{T}あたり \qty{42.58}{MHz}です。これを使うのが\textbf{核磁気共鳴}（NMR）です（ブロッホとパーセルが精密にはかる方法を作り、1952年のノーベル物理学賞を受けました）。\bkidx{かくじききょうめい}{核磁気共鳴}病院の \textbf{MRI} は、体を超伝導の磁石の強い磁場（1.5〜\qty{3}{T}：電-6）の中に入れ、体じゅうの水の陽子に電波を当てて共鳴させ、陽子が出し返す電波をはかって、体の断面をうつします。場所によって磁場を少しずつ変えておくと、共鳴の周波数が場所ごとに違うので、どこの陽子の電波かがわかるのです。

熱-15 の言い方では、磁場の中の陽子は「2つの準位をもつ物」です。上の数と下の数の比は $e^{-\Delta E/(k_{\mathrm B}T)}$ で、1.5 T・体温では、上と下の数の差は、全体のわずか100万分の5ほどです。それでも、体の中の水には陽子がとてもたくさんあるので、はかれる大きさの電波になります。

\Needspace{10\baselineskip}\bksubsection{7}{パウリの排他律：1つの席に1個まで}
\label{ph:原-7:sec7}
\textbf{1つの量子数の組 $(n,l,m_l,m_s)$ には、電子は1個しか入れない}。これを\textbf{パウリの排他律}といいます。\bkidx{ぱうりのはいたりつ}{パウリの排他律}パウリは、原子のスペクトルと周期表の性質を詳しく調べてこの決まりを見つけ、1945年のノーベル物理学賞を受けました。熱-15 ⑨では、「電子は、1つの状態（棚）に1個までしか入れない」と書きました。その「棚」が、ここでの席です。

\phboxhead{もし、排他律がなかったら}
排他律がなければ、どの原子でも、すべての電子がいちばん低い $1s$ の席に集まってしまいます。そうすると、原子番号が大きい原子ほど、電子の雲は原子核の近くに小さく縮み（③の2倍の図）、外側に「いちばん外の電子」というものがなくなります。元素の性質は周期的に繰り返さず、ナトリウムもネオンも塩素も、同じような性質になるでしょう。化学反応も、私たちの体も、ありえません。また、物が押してもつぶれず、決まった体積をもつのも、電子どうしが同じ席に入れず、雲どうしが重なりにくいためです。

\phboxhead{フェルミ粒子とボース粒子}
電子のように、1つの状態に1個までしか入れない粒子を\textbf{フェルミ粒子}といいます。陽子や中性子もフェルミ粒子です。反対に、光子のように、1つの状態にいくつでも入れる粒子を\textbf{ボース粒子}といいます（熱-15 ⑨のフェルミ分布とボース分布の違い）。原-8 では、金属の中のたくさんの電子を、この決まりで席に入れていきます。

\begin{column}{ボーアの円から、シュレディンガーの雲へ}
ボーアの模型（原-4）は、水素原子の光の色をみごとに説明し、ボーアは1922年のノーベル物理学賞を受けました。けれども、電子が2個のヘリウムになると、どうしても正しいエネルギーが出ませんでした。電子も波だ、というド・ブロイの考え（原-3。1929年のノーベル物理学賞）をもとに、シュレディンガーは波の方程式を作り（原-5。1933年のノーベル物理学賞）、水素原子では、ボーアと同じエネルギーと、円ではなく雲としての電子の姿を出しました。そこに、パウリの排他律（1945年のノーベル物理学賞）とスピンが加わって、周期表の 2・8・8・18 が、はじめて「なぜ」から説明できるようになったのです。メンデレーエフが表を作ってから、約60年後のことでした。
\end{column}

\Needspace{12\baselineskip}\bksubsection{8}{フントの規則と、磁石のもと}
\label{ph:原-7:sec8}
同じ $n$ と $l$ の席（たとえば $2p$ の3つの箱）に、電子が何個か入るとき、入り方にはもう1つの決まりがあります。\textbf{電子は、なるべく別々の箱に、スピンの向きをそろえて入る}。これを\textbf{フントの規則}といいます。\bkidx{ふんとのきそく}{フントの規則}
\begin{bkfigure}
\begin{tikzpicture}[x=1mm, y=1mm]
  \node[txh lab, anchor=south west] at (0,7) {C（$2p^2$）};
  \txhbox{(0,0)}{1}{0}\txhbox{(6,0)}{1}{0}\txhbox{(12,0)}{0}{0}
  \node[txh lab, anchor=south west] at (40,7) {N（$2p^3$）};
  \txhbox{(40,0)}{1}{0}\txhbox{(46,0)}{1}{0}\txhbox{(52,0)}{1}{0}
  \node[txh lab, anchor=south west] at (80,7) {O（$2p^4$）};
  \txhbox{(80,0)}{1}{1}\txhbox{(86,0)}{1}{0}\txhbox{(92,0)}{1}{0}
  \node[txh lab, anchor=south west] at (0,-11) {Fe（$3d^6$）};
  \txhbox{(0,-18)}{1}{1}\txhbox{(6,-18)}{1}{0}\txhbox{(12,-18)}{1}{0}\txhbox{(18,-18)}{1}{0}\txhbox{(24,-18)}{1}{0}
  \node[txh lab, anchor=north west] at (0,-19.5) {相手のいない電子 4個};
  \node[txh lab, anchor=south west] at (60,-11) {Cu（$3d^{10}$）};
  \txhbox{(60,-18)}{1}{1}\txhbox{(66,-18)}{1}{1}\txhbox{(72,-18)}{1}{1}\txhbox{(78,-18)}{1}{1}\txhbox{(84,-18)}{1}{1}
  \node[txh lab, anchor=north west] at (60,-19.5) {相手のいない電子 0個};
\end{tikzpicture}
\bkfigcaption{フントの規則による入り方。箱は $m_l$ の違う席で、上向きの矢印（左寄り）と下向きの矢印（右寄り）がスピンの向き。鉄の $3d$ には、向きのそろった、相手のいない電子が4個ある。銅の $3d$ は満席で、上向きと下向きが打ち消し合う。}\label{ph:原-7:fig-hund}
\end{bkfigure}

なぜでしょう。電子どうしは、電気の力で反発します。別々の箱（雲の向きが違う）に入れば、離れていられるので、反発のエネルギーが小さくなります。さらに、同じ向きのスピンの電子どうしは、パウリの排他律のために同じ場所に来ることができず、自然に避け合うので、反発がもっと小さくなります。この「スピンをそろえると、反発が減ってエネルギーが下がる」働きを\textbf{交換相互作用}といいます。\bkidx{こうかんそうごさよう}{交換相互作用}電-7 ⑩で「量子の決まりと電子どうしの電気の反発から生まれる」と予告した働きは、これです。磁石どうしの磁場の力ではなく、電気の力と排他律の組み合わせなので、ずっと強いのです。

\phboxhead{磁石になる原子・ならない原子}
鉄の原子の $3d$ には、向きのそろった電子が4個あり、原子1個が強い小さな磁石です。鉄の金属の中では、となりの原子どうしの電子の間にも交換相互作用が働き、原子の磁石の向きがそろいます。これが、磁区の中で向きがそろう強磁性のもとです（電-7 ⑩。金属の中の電子の帯の考えで、もっと詳しく見るのは原-8 ⑫）。銅の $3d$ は満席で、上向きと下向きが打ち消し合うので、原子は磁石になりません。ネオジム磁石のネオジムは、$4f$ の席に、向きのそろった電子を何個かもっています。

\Needspace{12\baselineskip}\bksubsection{9}{イオンになりやすさと、はてな？の答え}
\label{ph:原-7:sec9}
原子からいちばん外の電子を1個取り去るのに要るエネルギーが、電離エネルギーです（③）。図\ref{ph:原-7:fig-ie}は、原子番号1〜54 の原子の電離エネルギーを、はかった値で並べたものです。

\begin{bkfigure}
\begin{tikzpicture}
\begin{axis}[bkplot, axis lines=left, width=118mm, height=56mm, xmin=0, xmax=56, ymin=0, ymax=28, xtick={1,10,18,36,54}, ytick={0,5,10,15,20,25}, xlabel={原子番号 $Z$}, ylabel={電離エネルギー（eV）}]
  \draw[bkBlue!80!black, line width=1.0pt] (axis cs:1,13.60) -- (axis cs:2,24.59) -- (axis cs:3,5.39) -- (axis cs:4,9.32) -- (axis cs:5,8.30) -- (axis cs:6,11.26) -- (axis cs:7,14.53) -- (axis cs:8,13.62) -- (axis cs:9,17.42) -- (axis cs:10,21.56) -- (axis cs:11,5.14) -- (axis cs:12,7.65) -- (axis cs:13,5.99) -- (axis cs:14,8.15) -- (axis cs:15,10.49) -- (axis cs:16,10.36) -- (axis cs:17,12.97) -- (axis cs:18,15.76) -- (axis cs:19,4.34) -- (axis cs:20,6.11) -- (axis cs:21,6.56) -- (axis cs:22,6.83) -- (axis cs:23,6.75) -- (axis cs:24,6.77) -- (axis cs:25,7.43) -- (axis cs:26,7.90) -- (axis cs:27,7.88) -- (axis cs:28,7.64) -- (axis cs:29,7.73) -- (axis cs:30,9.39) -- (axis cs:31,6.00) -- (axis cs:32,7.90) -- (axis cs:33,9.79) -- (axis cs:34,9.75) -- (axis cs:35,11.81) -- (axis cs:36,14.00) -- (axis cs:37,4.18) -- (axis cs:38,5.69) -- (axis cs:39,6.22) -- (axis cs:40,6.63) -- (axis cs:41,6.76) -- (axis cs:42,7.09) -- (axis cs:43,7.12) -- (axis cs:44,7.36) -- (axis cs:45,7.46) -- (axis cs:46,8.34) -- (axis cs:47,7.58) -- (axis cs:48,8.99) -- (axis cs:49,5.79) -- (axis cs:50,7.34) -- (axis cs:51,8.61) -- (axis cs:52,9.01) -- (axis cs:53,10.45) -- (axis cs:54,12.13);
  \fill[bkBlue!80!black] (axis cs:1,13.60) circle[radius=1.0pt]; \fill[bkBlue!80!black] (axis cs:2,24.59) circle[radius=1.0pt]; \fill[bkBlue!80!black] (axis cs:3,5.39) circle[radius=1.0pt]; \fill[bkBlue!80!black] (axis cs:4,9.32) circle[radius=1.0pt]; \fill[bkBlue!80!black] (axis cs:5,8.30) circle[radius=1.0pt]; \fill[bkBlue!80!black] (axis cs:6,11.26) circle[radius=1.0pt]; \fill[bkBlue!80!black] (axis cs:7,14.53) circle[radius=1.0pt]; \fill[bkBlue!80!black] (axis cs:8,13.62) circle[radius=1.0pt]; \fill[bkBlue!80!black] (axis cs:9,17.42) circle[radius=1.0pt]; \fill[bkBlue!80!black] (axis cs:10,21.56) circle[radius=1.0pt]; \fill[bkBlue!80!black] (axis cs:11,5.14) circle[radius=1.0pt]; \fill[bkBlue!80!black] (axis cs:12,7.65) circle[radius=1.0pt]; \fill[bkBlue!80!black] (axis cs:13,5.99) circle[radius=1.0pt]; \fill[bkBlue!80!black] (axis cs:14,8.15) circle[radius=1.0pt]; \fill[bkBlue!80!black] (axis cs:15,10.49) circle[radius=1.0pt]; \fill[bkBlue!80!black] (axis cs:16,10.36) circle[radius=1.0pt]; \fill[bkBlue!80!black] (axis cs:17,12.97) circle[radius=1.0pt]; \fill[bkBlue!80!black] (axis cs:18,15.76) circle[radius=1.0pt]; \fill[bkBlue!80!black] (axis cs:19,4.34) circle[radius=1.0pt]; \fill[bkBlue!80!black] (axis cs:20,6.11) circle[radius=1.0pt]; \fill[bkBlue!80!black] (axis cs:21,6.56) circle[radius=1.0pt]; \fill[bkBlue!80!black] (axis cs:22,6.83) circle[radius=1.0pt]; \fill[bkBlue!80!black] (axis cs:23,6.75) circle[radius=1.0pt]; \fill[bkBlue!80!black] (axis cs:24,6.77) circle[radius=1.0pt]; \fill[bkBlue!80!black] (axis cs:25,7.43) circle[radius=1.0pt]; \fill[bkBlue!80!black] (axis cs:26,7.90) circle[radius=1.0pt]; \fill[bkBlue!80!black] (axis cs:27,7.88) circle[radius=1.0pt]; \fill[bkBlue!80!black] (axis cs:28,7.64) circle[radius=1.0pt]; \fill[bkBlue!80!black] (axis cs:29,7.73) circle[radius=1.0pt]; \fill[bkBlue!80!black] (axis cs:30,9.39) circle[radius=1.0pt]; \fill[bkBlue!80!black] (axis cs:31,6.00) circle[radius=1.0pt]; \fill[bkBlue!80!black] (axis cs:32,7.90) circle[radius=1.0pt]; \fill[bkBlue!80!black] (axis cs:33,9.79) circle[radius=1.0pt]; \fill[bkBlue!80!black] (axis cs:34,9.75) circle[radius=1.0pt]; \fill[bkBlue!80!black] (axis cs:35,11.81) circle[radius=1.0pt]; \fill[bkBlue!80!black] (axis cs:36,14.00) circle[radius=1.0pt]; \fill[bkBlue!80!black] (axis cs:37,4.18) circle[radius=1.0pt]; \fill[bkBlue!80!black] (axis cs:38,5.69) circle[radius=1.0pt]; \fill[bkBlue!80!black] (axis cs:39,6.22) circle[radius=1.0pt]; \fill[bkBlue!80!black] (axis cs:40,6.63) circle[radius=1.0pt]; \fill[bkBlue!80!black] (axis cs:41,6.76) circle[radius=1.0pt]; \fill[bkBlue!80!black] (axis cs:42,7.09) circle[radius=1.0pt]; \fill[bkBlue!80!black] (axis cs:43,7.12) circle[radius=1.0pt]; \fill[bkBlue!80!black] (axis cs:44,7.36) circle[radius=1.0pt]; \fill[bkBlue!80!black] (axis cs:45,7.46) circle[radius=1.0pt]; \fill[bkBlue!80!black] (axis cs:46,8.34) circle[radius=1.0pt]; \fill[bkBlue!80!black] (axis cs:47,7.58) circle[radius=1.0pt]; \fill[bkBlue!80!black] (axis cs:48,8.99) circle[radius=1.0pt]; \fill[bkBlue!80!black] (axis cs:49,5.79) circle[radius=1.0pt]; \fill[bkBlue!80!black] (axis cs:50,7.34) circle[radius=1.0pt]; \fill[bkBlue!80!black] (axis cs:51,8.61) circle[radius=1.0pt]; \fill[bkBlue!80!black] (axis cs:52,9.01) circle[radius=1.0pt]; \fill[bkBlue!80!black] (axis cs:53,10.45) circle[radius=1.0pt]; \fill[bkBlue!80!black] (axis cs:54,12.13) circle[radius=1.0pt];
  \node[font=\sffamily\footnotesize, anchor=south] at (axis cs:2,24.99) {He};
  \node[font=\sffamily\footnotesize, anchor=south] at (axis cs:10,21.96) {Ne};
  \node[font=\sffamily\footnotesize, anchor=south] at (axis cs:18,16.16) {Ar};
  \node[font=\sffamily\footnotesize, anchor=south] at (axis cs:36,14.40) {Kr};
  \node[font=\sffamily\footnotesize, anchor=south] at (axis cs:54,12.53) {Xe};
  \node[font=\sffamily\footnotesize, anchor=north] at (axis cs:3,4.89) {Li};
  \node[font=\sffamily\footnotesize, anchor=north] at (axis cs:11,4.64) {Na};
  \node[font=\sffamily\footnotesize, anchor=north] at (axis cs:19,3.84) {K};
  \node[font=\sffamily\footnotesize, anchor=north] at (axis cs:37,3.68) {Rb};
  \node[font=\sffamily\footnotesize, anchor=west] at (axis cs:2.2,13.6) {H};
\end{axis}
\end{tikzpicture}
\bkfigcaption{中性の原子の電離エネルギー（NIST の原子スペクトルのデータベースの値）。殻がうまった貴ガス（He・Ne・Ar・Kr・Xe）で山になり、そのすぐ次の、1個だけ外に電子をもつ原子（Li・Na・K・Rb）で、がくんと小さくなる。}\label{ph:原-7:fig-ie}
\end{bkfigure}

\phboxhead{有効核電荷：外の電子が感じる電荷}
ナトリウムの $3s$ の電子は、原子核の $+11e$ に引かれますが、そのあいだに内側の電子10個がいて、電荷の大部分を打ち消します（遮へい）。外の電子が感じる、遮へいのあとの電荷を $Z_{\mathrm{eff}}e$ と書き、$Z_{\mathrm{eff}}$ を\textbf{有効核電荷}とよびます。\bkidx{ゆうこうかくでんか}{有効核電荷}電子1個の原子の式（③の2倍の図）$E=-13.6\,Z^2/n^2$ eV の $Z$ を $Z_{\mathrm{eff}}$ に置きかえて、電離エネルギー $E_{\mathrm i}$ から逆に求めると、
\pheq{Z_{\mathrm{eff}}=n\sqrt{\frac{E_{\mathrm i}}{\qty{13.6}{eV}}}}{1 ＝ 1}
\label{ph:原-7:zeff}です。ナトリウムは $n=3$、$E_{\mathrm i}=\qty{5.14}{eV}$ で $Z_{\mathrm{eff}}\approx1.8$、ネオンは $n=2$、$E_{\mathrm i}=\qty{21.6}{eV}$ で $Z_{\mathrm{eff}}\approx2.5$ です（例題1）。完全に遮へいされるならナトリウムは1のはずですが、$3s$ の雲は内側にしみこむので、少し大きくなっています。

\phboxhead{イオン}
\label{ph:原-7:ion}原子が電子を失うと、電子の負の電荷が減るので、全体が正の\textbf{陽イオン}になります。電子を受け取ると、負の\textbf{陰イオン}になります（中学理科）。\bkidx{いおん}{イオン}ナトリウムは、外の電子1個を失うと、ネオンと同じ満席の電子配置の Na$^+$ になります。塩素（$3s^2\,3p^5$）は、$3p$ の空いた席に1個受け取ると、アルゴンと同じ満席の Cl$^-$ になります。食塩（NaCl）は、この2つのイオンが電気の力で引き合ってできた結晶です。周期表の左の原子（電離エネルギーが小さい）は陽イオンに、右の原子（貴ガスを除く。席があと少しでうまる）は陰イオンになりやすいのです。結合の中で電子を自分の側へ引き寄せる強さを\textbf{電気陰性度}といい、周期表の右上ほど大きく、フッ素がいちばん大きくなります。

\begin{motto}[4]
\textbf{ヘリウムの2個の電子は、どれだけ反発し合っている？}\par
ヘリウム原子から電子を1個取り去るには \qty{24.6}{eV}、残った He$^+$ からもう1個取り去るには \qty{54.4}{eV}（例題4）要ります。だから、ヘリウム原子の2個の電子をばらばらにするには、合わせて $24.6+54.4=\qty{79.0}{eV}$ が要ります。もし電子どうしが反発しなければ、2個とも He$^+$ と同じ \qty{-54.4}{eV} の席にいるはずで、合わせて \qty{108.8}{eV} が要るはずです。その差 \qty{29.8}{eV}が、2個の電子どうしの反発のエネルギーです。電子2個の電気の位置エネルギー $ke^2/d$（$ke^2=\qty{1.44}{eV.nm}$）がこの大きさになる距離 $d$ は、$1.44\div29.8\approx\qty{0.048}{nm}$ で、ヘリウムの電子の雲の大きさと同じくらいです。電子どうしの反発は、原子のエネルギーの中で、無視できない大きさなのです。
\end{motto}

\phboxhead{\bkhatenaref の答え}
\label{ph:原-7:hatena}ネオン（$1s^2\,2s^2\,2p^6$）は、$n=1$ の殻（2席）と $n=2$ の殻（8席）が、ちょうど満席です（④）。
\begin{enumerate}[label={\bkcircled{\arabic*}}, leftmargin=2.4\zw, itemsep=1pt, topsep=2pt]
\item \textbf{電子を手放さない}：いちばん外の $2p$ の電子は、$n=2$ と内側にいて、遮へいも弱く（$Z_{\mathrm{eff}}\approx2.5$）、取り去るのに \qty{21.6}{eV} も要ります。ナトリウムの外の電子（\qty{5.14}{eV}）の約4倍です。
\item \textbf{電子を受け取らない}：受け取った電子は、満席の $n=2$ には入れず、ずっと高い $3s$ の席に入るしかありません。そこでは原子核の電荷が10個の電子にほぼ打ち消されているので、電子を引き止められません。
\item \textbf{電子を分け合わない}：2つの原子が電子を分け合う結合（⑪）では、それぞれの原子の、相手のいない電子が、2つの原子のあいだの席に入ります。ネオンには、相手のいない電子が1個もありません。
\end{enumerate}
電子をやりとりしても分け合っても、エネルギーが下がらないので、ネオンはほかの物質とほとんど結びつかないのです。ナトリウムは反対に、満席の殻の外に電子が1個だけあり、それを簡単に手放してネオンと同じ形になれるので、とても反応しやすいのです。たった1個の違いが、「満席の外にいるかどうか」の違いになっています。

\phkokomade

\Needspace{9\baselineskip}\phstage{shiki}{シュレディンガー方程式で確かめる}{5}

\bksubsection{10}{1s を代入して、$-13.6$ eV を確かめる}
\label{ph:原-7:sec10}
原-5 ⑥の時間によらないシュレディンガー方程式を、3次元に広げます。2階の微分 $\dd^2\psi/\dd x^2$ を、3つの向きの2階の偏微分の和（ラプラシアン：8-12）に置きかえます。
\pheq{-\frac{\hbar^2}{2m_{\mathrm e}}\nabla^2\psi+U(r)\,\psi=E\,\psi,\qquad \nabla^2=\frac{\partial^2}{\partial x^2}+\frac{\partial^2}{\partial y^2}+\frac{\partial^2}{\partial z^2}}{J/m$^{3/2}$ ＝ J/m$^{3/2}$}
\label{ph:原-7:sch3d}水素原子では $U(r)=-ke^2/r$ です。$\psi$ が原子核からの距離 $r$ だけで決まる（丸い雲）とき、数学の本（8-12⑦）で作った球のラプラシアン
\[ \nabla^2\psi=\frac{1}{r^2}\,\frac{\dd}{\dd r}\Bigl(r^2\,\frac{\dd\psi}{\dd r}\Bigr) \]
が使えます。

\phboxhead{$\psi=e^{-r/a}$ を入れてみる}
$a$ はまだ決めていない長さです。
\begin{bkhenkei}[0.55]
\frac{\dd\psi}{\dd r}=-\frac{1}{a}\,e^{-r/a}=-\frac{\psi}{a} & $e^{cr}$ の微分（5-14⑤） \\
r^2\,\frac{\dd\psi}{\dd r}=-\frac{r^2}{a}\,\psi & $r^2$ をかける \\
\frac{\dd}{\dd r}\Bigl(-\frac{r^2}{a}\psi\Bigr)=-\frac{2r}{a}\psi+\frac{r^2}{a^2}\psi & 積の微分（5-6） \\
\nabla^2\psi=\Bigl(\frac{1}{a^2}-\frac{2}{ar}\Bigr)\psi & $r^2$ でわる \\
\end{bkhenkei}
これをシュレディンガー方程式の左辺に入れると、
\[ -\frac{\hbar^2}{2m_{\mathrm e}}\Bigl(\frac{1}{a^2}-\frac{2}{ar}\Bigr)\psi-\frac{ke^2}{r}\,\psi=\Bigl[-\frac{\hbar^2}{2m_{\mathrm e}a^2}+\Bigl(\frac{\hbar^2}{m_{\mathrm e}a}-ke^2\Bigr)\frac{1}{r}\Bigr]\psi \]
右辺は $E\psi$ で、$E$ は $r$ によらない定数です。左辺の $[\ ]$ の中が、どの $r$ でも同じ定数になるには、$1/r$ の係数が0でなければなりません。
\pheq{\frac{\hbar^2}{m_{\mathrm e}a}=ke^2\quad\text{より}\quad a=\frac{\hbar^2}{m_{\mathrm e}ke^2}=a_0\approx\qty{52.9}{pm}}{m ＝ (J·s)$^2$ ÷ (kg × J·m)}
そのとき、残った定数の部分がエネルギーです。
\pheq{E=-\frac{\hbar^2}{2m_{\mathrm e}a_0^{\,2}}=-\frac{ke^2}{2a_0}\approx\qty{-13.6}{eV}}{J ＝ (J·s)$^2$ ÷ (kg·m$^2$)}
\label{ph:原-7:E1}（2つめの等号は、$a_0$ の式から $\hbar^2/(m_{\mathrm e}a_0)=ke^2$ を使いました。）$e^{-r/a_0}$ は、確かにシュレディンガー方程式の解で、エネルギーは \qty{-13.6}{eV} です。原-4 のボーアの模型と同じ値が、円の軌道を考えずに出てきました。

\begin{phyomi}[$\hbar^2/(m_{\mathrm e}a)=ke^2$]
両辺を $a$ でわると、左辺の $\hbar^2/(m_{\mathrm e}a^2)$ は「雲を半径 $a$ に閉じこめたときの運動エネルギー」の目安（原-3 の不確定性：小さく閉じこめると運動量がばらつく）、右辺の $ke^2/a$ は「陽子が電子を引く位置エネルギー」の大きさです。雲を小さくすると、位置エネルギーは下がりますが、運動エネルギーが上がります。そのつり合う大きさが $a_0$ です。原子がつぶれずに、決まった大きさをもつのは、このつり合いのためです。
\end{phyomi}

\begin{phflow}[シュレディンガー方程式から、化学の式へ]
\phflowrow{$-\dfrac{\hbar^2}{2m_{\mathrm e}}\nabla^2\psi+U\psi=E\psi$}{$U=-ke^2/r$\\丸い雲（$r$ だけ）}{$\psi=e^{-r/a_0}$、$a_0=\dfrac{\hbar^2}{m_{\mathrm e}ke^2}$}
\phflowrow{$\psi=e^{-r/a_0}$ を代入}{$r$ によらない\\部分がエネルギー}{$E_1=-\dfrac{ke^2}{2a_0}=\qty{-13.6}{eV}$}
\phflowrow{量子数 $(n,l,m_l,m_s)$}{1つの席に\\電子は1個}{殻の席の数 $2n^2$（2・8・18）}
\end{phflow}

\begin{phbaisuru}[2倍]{ためしの長さ $a$}
\noindent\begin{minipage}[c]{0.45\linewidth}\raggedright
$a$ を $a_0$ の2倍にして $e^{-r/(2a_0)}$ を入れると、$1/r$ の係数が0になりません。$H\psi/\psi$（左辺÷$\psi$）は $r$ によって変わり、原子核の近くで大きく負になります。つまり、シュレディンガー方程式の解ではありません。$a=a_0$ のときだけ、$H\psi/\psi$ がどこでも \qty{-13.6}{eV} の一定の値になります。
\end{minipage}\hfill
\begin{minipage}[c]{0.52\linewidth}\centering
\begin{tikzpicture}
\begin{axis}[bkplot, axis lines=left, width=42mm, height=44mm, xmin=0, xmax=5, ymin=-60, ymax=5, xtick={0,1,2,3,4,5}, ytick={-60,-40,-20,-13.6,0}, yticklabels={$-60$,$-40$,$-20$,,0}, xlabel={$r$（$a_0$ を単位に）}, ylabel={$H\psi/\psi$（eV）}]
  \addplot[bkcurve, domain=0.2:4.9] {-13.6};
  \addplot[bkcurve2, dashed, domain=0.25:4.9, samples=80] {-3.4-13.6/x};
  \node[font=\sffamily\footnotesize, anchor=north] at (axis cs:3.0,-16.5) {$a=a_0$（実線）};
  \node[font=\sffamily\footnotesize, anchor=west] at (axis cs:0.9,-40) {$a=2a_0$（破線）};
\end{axis}
\end{tikzpicture}
\end{minipage}
\end{phbaisuru}

\phboxhead{$l$ と $m_l$ の意味（結果を読む）}
丸くない雲は、$\psi$ が $r$ のほかに角度 $\theta,\varphi$ にもよります。球座標でシュレディンガー方程式を解くと（専門の本で扱います）、解は「$r$ だけの部分」×「角度だけの部分」の積の形になり、次のことがわかります。
\begin{enumerate}[label={\bkcircled{\arabic*}}, leftmargin=2.4\zw, itemsep=1pt, topsep=2pt]
\item 角度の部分が、球面の上の定在波になれるのは、$l=0,1,2,\ldots$ と、それぞれ $m_l=-l,\ldots,l$ の $2l+1$ 通りのときだけ。$l$ は、雲が原子核の周りを回る勢い（角運動量）の大きさを表し、その大きさは $\sqrt{l(l+1)}\,\hbar$ です。$s$ の雲（$l=0$）は、回る勢いが0です。
\item $r$ の部分が、遠くで0になる定在波になれるのは、$n\ge l+1$ のときだけで、エネルギーは $n$ だけで決まる $-13.6/n^2$ eV。だから $l$ は $n-1$ まで。
\item 例：$2s$ は $\psi\propto\bigl(2-r/a_0\bigr)e^{-r/(2a_0)}$ で、$r=2a_0$ に節（電子が見つからない球面）があります。$2p$（たての向き）は $\psi\propto r\,e^{-r/(2a_0)}\cos\theta$ で、$\theta=90^\circ$ の面が節です（図\ref{ph:原-7:fig-cloud}）。
\end{enumerate}
水素原子では、同じ $n$ の $n^2$ 個の雲（$l$ と $m_l$ の違い）がみな同じエネルギーです。このように、違う状態が同じエネルギーになることを\textbf{縮退}といいます。電子が2個以上の原子では、⑤の遮へいで、$l$ の違う雲の縮退が解けます。

\phboxhead{プログラムで、電子の雲を点で打つ}
$1s$ の $|\psi|^2$ に比例する確率で、空間に点を打ちます（3-17 の、たくさん繰り返すと割合が確率に近づく考え）。1辺 $12a_0$ の箱の中に、でたらめに点を500万個ばらまき、それぞれの点を $e^{-2r/a_0}$ の確率で残します。
\begin{phcode}[Python]
import numpy as np
rng = np.random.default\_rng(1)
N = 5000000
xyz = rng.uniform(-6, 6, (N, 3))         \# a0 を単位に点をばらまく
r = np.sqrt((xyz**2).sum(axis=1))        \# 原子核からの距離
keep = rng.random(N) < np.exp(-2 * r)    \# |ψ|\textasciicircum2 に比例する確率で残す
rk = r[keep]
print(len(rk))                           \# 残った点の数
print(round(rk.mean(), 2))               \# 距離の平均（a0 を単位に）
print(round((rk < 1).mean(), 2))         \# a0 より内側の点の割合
print(round((rk < 3).mean(), 2))         \# 3a0 より内側の点の割合
\end{phcode}
\begin{phcodeout}[実行結果]
9073
1.5
0.33
0.94
\end{phcodeout}
残った約9000個の点の距離の平均は $1.5a_0$（例題3）、$a_0$ より内側の割合は 0.33 で、例題2 の正確な値 0.32 と、点の数が有限なぶんの誤差の範囲で合います。図\ref{ph:原-7:fig-cloud}は、同じやり方で、$2s$・$2p$ の $|\psi|^2$ でも点を打ったものです。

\Needspace{12\baselineskip}\bksubsection{11}{化学結合と、分子の回転}
\label{ph:原-7:sec11}
\phboxhead{2つの井戸に分かれた電子：水素分子イオン}
陽子2個と電子1個の、いちばん簡単な分子（水素分子イオン H$_2^+$）を考えます。電子にとっては、陽子 A と陽子 B の周りに、すりばちの井戸が2つ並んでいます。原-6 ⑪の「2つの井戸」と同じように、電子の波は、片方の井戸の $1s$ の波（$\psi_{\mathrm A}$・$\psi_{\mathrm B}$）を、同じ向きにたしたもの $\psi_{\mathrm A}+\psi_{\mathrm B}$ と、反対向きにたしたもの $\psi_{\mathrm A}-\psi_{\mathrm B}$ の2つに分かれます。

\begin{bkfigure}
\begin{tikzpicture}
\begin{axis}[bkplot, axis lines=left, name=pa, width=60mm, height=40mm, xmin=-4, xmax=4, ymin=-0.2, ymax=1.35, xtick={-1,1}, xticklabels={A,B}, ytick={0}, xlabel={}, ylabel={$\psi$}, title={\footnotesize 結合（同じ向きにたす）}]
  \addplot[bkcurve, domain=-3.9:3.9, samples=200] {exp(-abs(x+1))+exp(-abs(x-1))};
  \addplot[bkGray, densely dotted, line width=0.8pt, domain=-3.9:3.9, samples=200] {exp(-abs(x+1))};
  \addplot[bkGray, densely dotted, line width=0.8pt, domain=-3.9:3.9, samples=200] {exp(-abs(x-1))};
\end{axis}
\begin{axis}[bkplot, axis lines=left, at={(pa.east)}, anchor=west, xshift=10mm, width=60mm, height=40mm, xmin=-4, xmax=4, ymin=-1.15, ymax=1.15, xtick={-1,1}, xticklabels={A,B}, ytick={0}, xlabel={}, ylabel={$\psi$}, title={\footnotesize 反結合（反対向きにたす）}]
  \addplot[bkcurve2, dashed, domain=-3.9:3.9, samples=200] {exp(-abs(x+1))-exp(-abs(x-1))};
  \draw[bkGray, line width=0.4pt] (axis cs:-4,0) -- (axis cs:4,0);
\end{axis}
\end{tikzpicture}
\bkfigcaption{2つの陽子 A・B を通る線の上の、電子の波。点線はそれぞれの陽子の $1s$ の波。左（実線）：同じ向きにたすと、2つの陽子のあいだにも電子が多くいる（結合）。右（破線）：反対向きにたすと、まん中で0になる節ができ、陽子のあいだに電子がいない（反結合）。}\label{ph:原-7:fig-bond}
\end{bkfigure}

$\psi_{\mathrm A}+\psi_{\mathrm B}$ では、電子が2つの陽子の\textbf{あいだ}にも多くいて、両方の陽子を引きつけ、陽子どうしの反発に勝って2つをつなぎとめます。エネルギーは、ばらばらの原子より低くなります。これが\textbf{結合}の状態です。$\psi_{\mathrm A}-\psi_{\mathrm B}$ では、まん中に節があり、陽子のあいだに電子がいないので、エネルギーは高く、2つの陽子は離れてしまいます（\textbf{反結合}）。\bkidx{きょうゆうけつごう}{共有結合}水素分子 H$_2$ では、2個の電子が、スピンを反対にして（パウリの排他律）、2個とも結合の席に入ります。こうして2つの原子が電子を分け合うのが\textbf{共有結合}です。H$_2$ をばらばらにするには \qty{4.48}{eV}（1 mol で約 \qty{432}{kJ}）が要ります（例題10）。

電-1 で「原子や分子のあいだの力のほとんどは、もとをたどるとクーロンの力だが、計算には量子力学が要る」と書いたのは、このことです。力は電気の力だけなのに、電子が波で、席に決まりがあるために、結合ができたりできなかったりします。2つの原子の距離を変えたときのエネルギーは、ある距離（H$_2$ では \qty{74}{pm}）で谷の底になります。谷の底の近くでは、エネルギーは距離のずれの2乗に比例し、原子どうしは「ばね」でつながれたようにふるまいます。力-8 で、ヤング率のもとと書いた原子のばねは、この谷の形です（ばね定数と分子の振動は原-6 ⑨）。

\phboxhead{分子の回転のエネルギーも、とびとび}
2つの原子からなる分子が、まん中を中心に回るとき、回転のエネルギーもとびとびになり、
\pheq{E_J=\frac{\hbar^2}{2I}\,J(J+1)\qquad(J=0,1,2,\ldots)}{J ＝ (J·s)$^2$ ÷ (kg·m$^2$)}
\label{ph:原-7:rot}です（結果を読みます）。$I=\mu r^2$ は分子の慣性モーメントで、$\mu=m_1m_2/(m_1+m_2)$ は換算質量（ボーア磁子 $\mu_{\mathrm B}$ とは別の量で、⑪だけで使います）、$r$ は原子の間の距離です。一酸化炭素 CO では $\hbar^2/(2I)$ が約 \qty{2.4e-4}{eV}で、$J=0$ から1へ移る差 $\hbar^2/I$ は、周波数 約 \qty{116}{GHz}の電波にあたります（例題11。観測される線は \qty{115.27}{GHz}）。宇宙の冷たいガスの雲の中の CO は、この電波を出していて、電波望遠鏡で星が生まれる場所を探すのに使われます。

段の間かく $\hbar^2/(2I)$ は、室温の $k_{\mathrm B}T$（0.026 eV）の約100分の1しかありません。だから、室温では回転のエネルギーの段は細かすぎて、ほとんど連続とみなせます。熱-14 で、分子の回転が等分配の法則どおりにエネルギーを受けとる（比熱に $\tfrac12k_{\mathrm B}$ ずつ効く）とした理由はここにあります。一方、分子の振動の段は 0.1〜0.5 eV と $k_{\mathrm B}T$ よりずっと大きく、室温では「こおって」います（原-6 ⑩）。

\Needspace{8\baselineskip}\bksubsection{12}{周期表は、どこまで続く？}
\label{ph:原-7:sec12}
天然にたくさんある元素は、ウラン（92番）くらいまでです。それより重い元素は、加速器で原子核どうしをぶつけて作ります。いまの周期表には118番のオガネソンまでが並び、第7周期がちょうどうまりました。113番の\textbf{ニホニウム}（Nh）は、理化学研究所の研究グループが、亜鉛の原子核をビスマスの原子核にぶつけて合成しました（2004年・2005年・2012年に合わせて3個）。2016年に名前が決まり、日本で見つかった初めての元素として周期表にのっています。こうした重い原子核は、1秒より短い時間でこわれてしまうものがほとんどです（原子核がこわれる話は原-12）。

\begin{phmodel}[電子1個の水素原子の解]
\phignore{陽子の動き（陽子は電子の約1840倍重いので、止まっているとした）。陽子の大きさ。相対論の効果と、スピンと軌道の磁気の力（とても小さなエネルギーのずれ）。}
\phassume{電子は1個で、陽子の電気の力 $-ke^2/r$ だけを受ける。}
\phlimit{電子が2個以上の原子（電子どうしの反発と遮へいで、$l$ によってエネルギーが違ってくる：⑤）。原子番号が大きく、内側の電子がとても速い原子（相対論が要る）。分子や固体の中の原子（となりの原子の電子と席を分け合う：⑪・原-8）。}
\end{phmodel}

\begin{phmada}[周期表は、どこまで続く？]
118番より重い元素（第8周期）を作る実験が、世界のいくつかの加速器で始まっていますが、まだ見つかっていません。原子番号が大きいほど、作れる数が少なくなり、すぐにこわれるからです。原子核にも「満席」で安定になる陽子・中性子の数があり、計算では、陽子114個・中性子184個のあたりに、ほかより長く生きる原子核の集まり（\textbf{安定の島}）があると予想されています。けれども、いま作れるいちばん中性子の多い114番の原子核でも、中性子が184個まで9個ほど足りず、島にはまだ届いていません。島があるのか、あるならどれくらい長く生きるのかは、わかっていません。

電子が多い原子や分子では、電子どうしが互いに影響し合うので、シュレディンガー方程式をそのまま正確に解くことは、コンピュータでもできません。量子化学の計算では、近似をくふうして、実験に合う答えを出しています。どこまで近似を減らせるかは、いまも研究が続いています。
\end{phmada}

\begin{phsaki}{原子から、分子・物質の設計へ}
原子や分子の中の電子の状態を、コンピュータで計算する分野を\textbf{量子化学}・\textbf{第一原理計算}といいます。電池の材料、薬の分子、触媒、半導体の新しい材料を、作る前に計算で調べることができます。分子の振動・回転のスペクトルから物の成分をはかる分光学、磁性の物理（スピントロニクス）、MRI の医用工学も、この章の続きです。詳しくは、量子化学・物性物理・磁性材料・医用画像の専門の本で学びます。
\end{phsaki}

\bkkigocards{\bkkigocard{$l$}{エル}{方位量子数。雲の形を決める数。長さ $l$ とは別}{azimuthal quantum number}{原-7}}
{\bkkigocard{$\mu_{\mathrm B}$}{ミュー・ビー}{ボーア磁子〔J/T〕。$e\hbar/(2m_{\mathrm e})$。電子の磁石の強さの目もり}{Bohr magneton}{原-7}}

\begin{taisetsu}
\def\tsl#1{\makebox[7.2\zw][l]{\sffamily\bfseries\small\color{\bkTT}#1}}
\tsl{水素原子}$E_n=-\dfrac{ke^2}{2a_0}\dfrac{1}{n^2}=-\dfrac{\qty{13.6}{eV}}{n^2}$、$a_0=\dfrac{\hbar^2}{m_{\mathrm e}ke^2}=\qty{52.9}{pm}$\\[3pt]
\tsl{1s の雲}$\psi\propto e^{-r/a_0}$、動径分布 $P(r)\propto r^2e^{-2r/a_0}$ の山は $r=a_0$\\[3pt]
\tsl{量子数}$n=1,2,\ldots$、$l=0,\ldots,n-1$、$m_l=-l,\ldots,l$、$m_s=\pm\tfrac12$\\[3pt]
\tsl{席の数}$l$ ごとに $2(2l+1)$（2・6・10・14）、殻ごとに $2n^2$（2・8・18・32）\\[3pt]
\tsl{スピン}$\Delta E=g\mu_{\mathrm B}B$、ESR 28.0 GHz/T、陽子の NMR 42.58 MHz/T\\[3pt]
\tsl{有効核電荷}$Z_{\mathrm{eff}}=n\sqrt{E_{\mathrm i}/\qty{13.6}{eV}}$
\end{taisetsu}

\begin{senmonsho}[量子力学・量子化学の専門書の書き方]
\begin{bkhonbun}
水素原子のハミルトニアンは
\[ \hat H=-\frac{\hbar^2}{2m}\nabla^2-\frac{e^2}{4\pi\varepsilon_0r} \]
である．固有関数は $\psi_{nlm}(r,\theta,\phi)=R_{nl}(r)Y_l^m(\theta,\phi)$ と変数分離され，エネルギー固有値は
\[ E_n=-\frac{me^4}{2(4\pi\varepsilon_0)^2\hbar^2n^2} \]
となり，$l,m$ について $n^2$ 重に縮退している．
\end{bkhonbun}
\bkyaku{「ハミルトニアン」$\hat H$ は、この章のシュレディンガー方程式の左辺の、$\psi$ にかける部分全体を1つの記号で書いたもの（$\hat H\psi=E\psi$。記号の上の山形は「演算子」の印：原-5）。$\frac{1}{4\pi\varepsilon_0}$ はこの本の $k$。$Y_l^m$ は「球面調和関数」で、角度の部分の定在波。$E_n$ の式に $k$ と $a_0$ を入れると、この章の $-ke^2/(2a_0n^2)$ と同じになる。}
\end{senmonsho}

\begin{kotoba}
\small
\bkword{量子数}{原子の中の電子の状態を区別する4つの数 $n$・$l$・$m_l$・$m_s$。}
\bkword{パウリの排他律}{1つの量子数の組（席）には、電子は1個しか入れない。}
\bkword{電子配置}{どの席に何個の電子が入っているか。$1s^2\,2s^2\,2p^6\,3s^1$ のように書く。}
\bkword{スピン}{電子が生まれつきもつ、角運動量と小さな磁石の性質。向きは2つ。}
\bkword{フントの規則}{同じ $n$・$l$ の席には、なるべく別々の箱に、スピンをそろえて入る。}
\bkword{遮へい・有効核電荷}{内側の電子が原子核の電荷を打ち消すこと。外の電子が感じる電荷の数が有効核電荷。}
\bkword{電離エネルギー}{原子からいちばん外の電子を1個取り去るのに要るエネルギー。}
\bkword{共有結合}{2つの原子が、相手のいない電子を1個ずつ出し合い、あいだの席で分け合う結合。}
\end{kotoba}

\bkkigocards{\bkkigocard{$Z_{\mathrm{eff}}$}{ゼット・エフ}{有効核電荷。外の電子が感じる、遮へいのあとの原子核の電荷の数}{effective nuclear charge}{原-7}}{}

%% ======================== 例題 ========================
\Needspace{10\baselineskip}\bksection{reidai}{例題でたしかめよう}
\label{ph:原-7:ex1}
\begin{phreidai}[ネオンとナトリウム（はてな？）]
ネオン（$Z=10$）の電離エネルギーは \qty{21.6}{eV}、ナトリウム（$Z=11$）は \qty{5.14}{eV}です（NIST の値）。いちばん外の電子の $n$ を電子配置から決め、有効核電荷 $Z_{\mathrm{eff}}$ を求めましょう。また、遮へいがまったくないとしたら、電離エネルギーはいくらになるはずですか。
\phbamen
貴ガスと、そのすぐ次の原子の、電子の離れやすさを比べる場面です。
\phmushi
内側の電子の影響は、原子核の電荷を $Z_{\mathrm{eff}}e$ に減らす働き（遮へい）にまとめ、電子1個の原子の式を使います。
\phzu
ネオン：$1s^2\,2s^2\,2p^6$、外の電子は $n=2$。ナトリウム：$1s^2\,2s^2\,2p^6\,3s^1$、外の電子は $n=3$。
\phriyu
電子1個の原子のエネルギー $-13.6\,Z^2/n^2$ eV（③）の $Z$ を $Z_{\mathrm{eff}}$ とし、その大きさが電離エネルギーだと考えます（⑨）。
\phtateru
\[ Z_{\mathrm{eff}}=n\sqrt{\frac{E_{\mathrm i}}{13.6}},\qquad \text{遮へいなし：}\ E=13.6\times\frac{Z^2}{n^2} \]
\phtoku
\[ \text{ネオン：}Z_{\mathrm{eff}}=2\times\sqrt{\frac{21.6}{13.6}}\approx2.5,\qquad \text{ナトリウム：}Z_{\mathrm{eff}}=3\times\sqrt{\frac{5.14}{13.6}}\approx1.8 \]
\[ \text{遮へいなし：ネオン}\ 13.6\times\frac{10^2}{2^2}=\qty{340}{eV},\qquad \text{ナトリウム}\ 13.6\times\frac{11^2}{3^2}\approx\qty{183}{eV} \]
\phtashikame
単位：根号の中は eV÷eV で、$Z_{\mathrm{eff}}$ は単位のない数。\\
大きさ：どちらも1と $Z$ のあいだ。ナトリウムは、内側の10個にほとんど打ち消されて2より小さい。ネオンは、同じ $n=2$ の電子どうしではあまり打ち消し合わないので、大きい。\\
極端な場合：遮へいがなければ、電離エネルギーは何百 eV にもなり、実際の値（数 eV〜20 eV ほど）と合いません。遮へいはとても大きいのです。
\kotae ネオン $Z_{\mathrm{eff}}\approx2.5$、ナトリウム $Z_{\mathrm{eff}}\approx1.8$。遮へいなしなら \qty{340}{eV} と \qty{183}{eV}
\end{phreidai}
\begin{jibun}
アルゴン（$Z=18$、$1s^2\,2s^2\,2p^6\,3s^2\,3p^6$）の電離エネルギーは \qty{15.8}{eV}です。いちばん外の電子の $n$ と、有効核電荷を求めましょう。
\bkjans{$n=3$。$Z_{\mathrm{eff}}=3\sqrt{15.8/13.6}\approx3.2$}
\end{jibun}

\label{ph:原-7:ex2}
\begin{phreidai}[1s の電子が $a_0$ の内側にいる確率]
水素原子の $1s$ の電子が、原子核から $a_0$ より内側で見つかる確率を求めましょう。$F(x)=-e^{-x}(x^2+2x+2)$ を微分すると $x^2e^{-x}$ になることを使ってかまいません。
\phbamen
雲のどこまでに、電子がどれだけの確率でいるかを調べる場面です（はてな？の図の円の内側）。
\phmushi
陽子は止まっている点とします。
\phzu
動径分布 $P(r)=\dfrac{4r^2}{a_0^{\,3}}e^{-2r/a_0}$（③）。求める確率は、$r=0$ から $a_0$ までの $P(r)$ の面積。
\phriyu
$P(r)\,\Delta r$ が、うすい殻にいる確率なので、それを $0$ から $a_0$ まで合わせる（積分する：5-10）と、内側にいる確率になります。
\phtateru
$x=2r/a_0$ とおくと、$r=0$ で $x=0$、$r=a_0$ で $x=2$、$\dd r=\tfrac{a_0}{2}\dd x$ なので（置換積分：5-16）、
\[ \int_0^{a_0}\frac{4r^2}{a_0^{\,3}}e^{-2r/a_0}\,\dd r=\int_0^{2}\frac{4}{a_0^{\,3}}\Bigl(\frac{a_0x}{2}\Bigr)^2e^{-x}\,\frac{a_0}{2}\,\dd x=\frac12\int_0^2x^2e^{-x}\,\dd x \]
\phtoku
$F'(x)=x^2e^{-x}$ なので（自分で微分して確かめましょう）、
\[ \frac12\Bigl[F(x)\Bigr]_0^2=\frac12\bigl\{-e^{-2}(4+4+2)+2\bigr\}=1-5e^{-2}\approx0.32 \]
\phtashikame
単位：確率なので単位はない。\\
大きさ：いちばん見つかりやすい距離が $a_0$ でも、内側にいる確率は3分の1ほど。残りの約3分の2は、$a_0$ より外にいます。\\
極端な場合：上の限りを $\infty$ にすると、$e^{-x}(x^2+2x+2)\to0$ なので $\tfrac12\times2=1$。全部の空間では確率1になり、規格化と合います。
\kotae 約 0.32（$1-5e^{-2}$）
\end{phreidai}
\begin{jibun}
例題2 と同じやり方で、$1s$ の電子が $2a_0$ より内側で見つかる確率を求めましょう。
\bkjans{$\frac12\bigl[F(x)\bigr]_0^4=1-13e^{-4}\approx0.76$}
\end{jibun}

\label{ph:原-7:ex3}
\begin{phreidai}[1s の電子の平均の距離と、いちばん多い距離]
水素原子の $1s$ の電子について、原子核からの距離の平均 $\langle r\rangle$ と、いちばん見つかりやすい距離を、pm の単位で求めましょう。$\displaystyle\int_0^\infty x^3e^{-x}\,\dd x=6$ を使ってかまいません。
\phbamen
雲の「大きさ」を、2通りの数で表す場面です。
\phmushi
陽子は止まっている点とします。
\phzu
$P(r)=\dfrac{4r^2}{a_0^{\,3}}e^{-2r/a_0}$、$a_0=\qty{52.9}{pm}$。
\phriyu
平均は「距離×その距離にいる確率」の合計（期待値：3-17④）。いちばん多い距離は $P(r)$ の山で、③のもっと知りたいで求めた $r=a_0$。
\phtateru
\[ \langle r\rangle=\int_0^\infty r\,P(r)\,\dd r=\int_0^\infty\frac{4r^3}{a_0^{\,3}}e^{-2r/a_0}\,\dd r \]
\phtoku
$x=2r/a_0$ とおくと、$r^3=\tfrac{a_0^3x^3}{8}$、$\dd r=\tfrac{a_0}{2}\dd x$ なので、
\[ \langle r\rangle=\frac{4}{a_0^{\,3}}\cdot\frac{a_0^{\,3}}{8}\cdot\frac{a_0}{2}\int_0^\infty x^3e^{-x}\,\dd x=\frac{a_0}{4}\times6=1.5a_0\approx\qty{79}{pm} \]
いちばん見つかりやすい距離は $a_0\approx\qty{53}{pm}$。
\phtashikame
単位：$a_0$ と同じ長さ。\\
大きさ：平均は、いちばん多い距離より大きい。$P(r)$ の山は、外の側にすそが長く引いているからです（図\ref{ph:原-7:fig-radial}）。\\
極端な場合：⑩のプログラムで打った点の距離の平均も 1.5 でした。
\kotae 平均 約 \qty{79}{pm}（$1.5a_0$）、いちばん見つかりやすいのは 約 \qty{53}{pm}（$a_0$）
\end{phreidai}
\begin{jibun}
電子1個のヘリウムイオン He$^+$（$Z=2$）では、$1s$ の動径分布は $P(r)\propto r^2e^{-4r/a_0}$ です。いちばん見つかりやすい距離と、距離の平均を求めましょう。
\bkjans{山は $r=a_0/2\approx\qty{26}{pm}$。平均は $0.75a_0\approx\qty{40}{pm}$（どちらも水素の半分）}
\end{jibun}

\label{ph:原-7:ex4}
\begin{phreidai}[電子1個のヘリウムイオンを、代入で解く]
電子1個のヘリウムイオン He$^+$ では、原子核の電荷が $+2e$ で、$U=-2ke^2/r$ です。$\psi=e^{-r/a}$ をシュレディンガー方程式に代入して、$a$ とエネルギーを求めましょう。
\phbamen
⑩の方法を、原子核の電荷の違う原子に使う場面です。
\phmushi
原子核は止まっている点とします。電子は1個なので、電子どうしの反発はありません。
\phzu
$\nabla^2\psi=\Bigl(\dfrac{1}{a^2}-\dfrac{2}{ar}\Bigr)\psi$（⑩で計算済み）、$U=-\dfrac{2ke^2}{r}$。
\phriyu
$\psi$ が $r$ だけで決まるので、球のラプラシアン（8-12⑦）で計算でき、⑩と同じ形の式になります。
\phtateru
\[ -\frac{\hbar^2}{2m_{\mathrm e}}\Bigl(\frac{1}{a^2}-\frac{2}{ar}\Bigr)\psi-\frac{2ke^2}{r}\psi=\Bigl[-\frac{\hbar^2}{2m_{\mathrm e}a^2}+\Bigl(\frac{\hbar^2}{m_{\mathrm e}a}-2ke^2\Bigr)\frac1r\Bigr]\psi=E\psi \]
\phtoku
$1/r$ の係数が0：$\dfrac{\hbar^2}{m_{\mathrm e}a}=2ke^2$ より、$a=\dfrac{\hbar^2}{2m_{\mathrm e}ke^2}=\dfrac{a_0}{2}\approx\qty{26.5}{pm}$。
\[ E=-\frac{\hbar^2}{2m_{\mathrm e}a^2}=-\frac{\hbar^2}{2m_{\mathrm e}}\cdot\frac{4}{a_0^{\,2}}=4\times(\qty{-13.6}{eV})\approx\qty{-54.4}{eV} \]
\phtashikame
単位：⑩と同じ。\\
大きさ：③の2倍の図のとおり、雲は半分、エネルギーは4倍の深さ。$-13.6\,Z^2$ eV の形に合います。\\
極端な場合：$Z=1$ にもどすと、⑩の水素の答えになります。
\kotae $a=a_0/2\approx\qty{26.5}{pm}$、$E\approx\qty{-54.4}{eV}$
\end{phreidai}
\begin{jibun}
電子1個のリチウムイオン Li$^{2+}$（$Z=3$）で、同じように $a$ とエネルギーを求めましょう。
\bkjans{$a=a_0/3\approx\qty{17.6}{pm}$、$E=9\times(-13.6)\approx\qty{-122}{eV}$}
\end{jibun}

\label{ph:原-7:ex5}
\begin{phreidai}[リチウムの 2s の電子が感じる電荷]
リチウム（$Z=3$、$1s^2\,2s^1$）の電離エネルギーは \qty{5.39}{eV}です。$2s$ の電子が感じる有効核電荷を求め、内側の $1s$ の電子2個が、原子核の電荷をどれだけ打ち消しているかを見積もりましょう。
\phbamen
いちばん簡単な「内側の殻と外の電子」の組で、遮へいの大きさを数にする場面です。
\phmushi
内側の電子の影響を、有効核電荷にまとめます。
\phzu
$n=2$、$E_{\mathrm i}=\qty{5.39}{eV}$、$Z=3$。
\phriyu
有効核電荷の式（⑨）を使います。
\phtateru
\[ Z_{\mathrm{eff}}=n\sqrt{\frac{E_{\mathrm i}}{13.6}},\qquad \text{打ち消した分}=Z-Z_{\mathrm{eff}} \]
\phtoku
\[ Z_{\mathrm{eff}}=2\times\sqrt{\frac{5.39}{13.6}}\approx1.26,\qquad 3-1.26\approx1.7 \]
\phtashikame
単位：単位のない数。\\
大きさ：内側の電子は2個なので、全部打ち消せば $Z-2=1$。実際は1より少し大きく、$2s$ の雲が内側に少ししみこんで、遮へいされない時間があることを表しています。\\
極端な場合：$2s$ の電子がずっと外にいるなら、$Z_{\mathrm{eff}}$ は1に近づきます。
\kotae $Z_{\mathrm{eff}}\approx1.26$。内側の2個が、約1.7個分の電荷を打ち消している
\end{phreidai}
\begin{jibun}
カリウム（$Z=19$、外の電子は $4s^1$）の電離エネルギーは \qty{4.34}{eV}です。有効核電荷を求めましょう。
\bkjans{$Z_{\mathrm{eff}}=4\sqrt{4.34/13.6}\approx2.3$}
\end{jibun}

\label{ph:原-7:ex6}
\begin{phreidai}[1 T の磁場の中の電子：ESR]
\qty{1.0}{T}の磁場の中で、電子のスピンの上と下のエネルギーの差と、ESR の周波数を求めましょう。また、\qty{300}{K}で、上の向き（エネルギーの高いほう）の数と下の向きの数の比を求めましょう。$g=2.00$、$\mu_{\mathrm B}=\qty{5.79e-5}{eV/T}$、$h=\qty{4.14e-15}{eV.s}$ とします。
\phbamen
熱-15 で「2つの準位をもつ物」とよんだものの、中身を求める場面です。
\phmushi
電子は1個ずつ別々で、となりの電子の影響は考えません。
\phzu
$\Delta E=g\mu_{\mathrm B}B$、$B=\qty{1.0}{T}$、$k_{\mathrm B}T=\qty{0.0259}{eV}$。
\phriyu
磁場の中のスピンのエネルギーの差（⑥）、光子のエネルギーと周波数 $hf=\Delta E$（原-2）、上と下の数の比はボルツマン因子（熱-15 ⑦）。
\phtateru
\[ \Delta E=g\mu_{\mathrm B}B,\qquad f=\frac{\Delta E}{h},\qquad \frac{\text{上の数}}{\text{下の数}}=e^{-\Delta E/(k_{\mathrm B}T)} \]
\phtoku
\[ \Delta E=2.00\times5.79\times10^{-5}\times1.0\approx\qty{1.16e-4}{eV} \]
\[ f=\frac{1.16\times10^{-4}}{4.14\times10^{-15}}\approx\qty{2.80e10}{Hz}=\qty{28.0}{GHz} \]
\[ \frac{\text{上の数}}{\text{下の数}}=e^{-1.16\times10^{-4}/0.0259}=e^{-0.00448}\approx0.9955 \]
\phtashikame
単位：(eV/T)×T＝eV、eV÷(eV·s)＝1/s＝Hz。\\
大きさ：28 GHz は、電子レンジ（2.45 GHz）の約10倍の周波数の電波です。上と下の数は、ほとんど同じで、差は0.5\,\%ほど。熱-15 の例（差 0.1 eV で比 0.021）に比べて $\Delta E$ がずっと小さいからです。\\
極端な場合：温度を下げるほど比は小さくなり、0 K ではすべてが下の向きにそろいます。
\kotae $\Delta E\approx\qty{1.16e-4}{eV}$、ESR の周波数 約 \qty{28.0}{GHz}、比 約 0.9955
\end{phreidai}
\begin{jibun}
多くの ESR の装置は、周波数 \qty{9.5}{GHz}の電波を使います。電子が共鳴する磁場の強さを求めましょう。
\bkjans{$B=hf/(g\mu_{\mathrm B})$ に数を入れて、約 \qty{0.34}{T}}
\end{jibun}

\label{ph:原-7:ex7}
\begin{phreidai}[MRI の周波数（3けたで書く）]
陽子の NMR の周波数は、\qty{1}{T}あたり \qty{42.58}{MHz}です（CODATA の値。3けたで書きます）。1.5 T と 3.0 T の MRI で、陽子が共鳴する電波の周波数を求めましょう。
\phbamen
病院の MRI の電波の周波数を決める場面です（電-6 の MRI の磁場）。
\phmushi
陽子の周りの電子が磁場を少し弱める効果（化学の環境によるずれ。100万分の数ほど）は無視します。
\phzu
$f=\gamma B$、$\gamma=\qty{42.58}{MHz/T}$。
\phriyu
共鳴の周波数は $\Delta E/h$ で、$\Delta E$ は $B$ に比例するので、周波数も $B$ に比例します（⑥）。
\phtateru
\[ f=42.58\times B \]
\phtoku
\[ f_{1.5}=42.58\times1.5\approx\qty{63.9}{MHz},\qquad f_{3.0}=42.58\times3.0\approx\qty{127.7}{MHz} \]
\phtashikame
単位：(MHz/T)×T＝MHz。\\
大きさ：FM ラジオ（76〜95 MHz）と同じくらいの周波数の電波です。電子の ESR（例題6）の約660分の1。\\
極端な場合：磁場が2倍なら周波数も2倍で、ちょうど2倍になっています。
\kotae 1.5 T で 約 \qty{63.9}{MHz}、3.0 T で 約 \qty{127.7}{MHz}
\end{phreidai}
\begin{jibun}
磁場 \qty{0.50}{T}の、小さな MRI の装置で、陽子が共鳴する周波数を求めましょう（3けたで書きます）。
\bkjans{$42.58\times0.50\approx\qty{21.3}{MHz}$}
\end{jibun}

\label{ph:原-7:ex8}
\begin{phreidai}[銅の電子配置から、自由電子の数を求める]
銅（$Z=29$）の電子配置 $[\mathrm{Ar}]\,3d^{10}\,4s^1$ から、金属の銅の中を動き回る電子（自由電子）が、原子1個あたり何個かを考え、1 m$^3$ あたりの自由電子の数を求めましょう。銅の密度は \qty{8.96}{g/cm^3}、モル質量は \qty{63.55}{g/mol}です。
\phbamen
電-4 で「自由電子の数」として使った数のもとを、電子配置から説明する場面です。
\phmushi
$3d$ の電子は原子に強く結びついていて動き回らないとし、$4s$ の電子だけが自由になると考えます。
\phzu
$3d^{10}$ は満席、$4s$ に1個。密度 $\rho=\qty{8.96e6}{g/m^3}$、モル質量 $M_{\text{モル}}=\qty{63.55}{g/mol}$、$N_{\mathrm A}=\num{6.02e23}\,\mathrm{/mol}$。1 m$^3$ あたりの自由電子の数を $n_{\text{電子}}$ と書きます（主量子数 $n$ と分けるため）。
\phriyu
いちばん外の、満席でない殻の電子は、となりの原子の席にも移れるので、金属の中を動き回れます（原-8 ④）。原子1個から1個です。
\phtateru
\[ n_{\text{電子}}=\frac{\rho}{M_{\text{モル}}}\times N_{\mathrm A}\times1 \]
\phtoku
\[ n_{\text{電子}}=\frac{8.96\times10^6}{63.55}\times6.02\times10^{23}\approx8.49\times10^{28}\ \text{個/m}^3 \]
\phtashikame
単位：(g/m$^3$)÷(g/mol)×(1/mol)＝1/m$^3$。\\
大きさ：電-4 で使った $8.49\times10^{28}$ と同じです。1 cm$^3$ に $8.5\times10^{22}$ 個。\\
極端な場合：$3d$ の電子も動き回るなら、数は11倍になるはずですが、電気の通しやすさや、原-8 で求めるフェルミエネルギーは、1個として計算した値によく合います。
\kotae 原子1個あたり1個、約 $8.49\times10^{28}$ 個/m$^3$
\end{phreidai}
\begin{jibun}
銀（$Z=47$、$[\mathrm{Kr}]\,4d^{10}\,5s^1$）の密度は \qty{10.50}{g/cm^3}、モル質量は \qty{107.9}{g/mol}です。1 m$^3$ あたりの自由電子の数を求めましょう。
\bkjans{原子1個に1個。$n_{\text{電子}}=\rho N_{\mathrm A}/M_{\text{モル}}$ に入れて、約 $5.86\times10^{28}$ 個/m$^3$}
\end{jibun}

\label{ph:原-7:ex9}
\begin{phreidai}[鉄の原子の磁石を全部そろえると]
鉄の金属の、原子1個あたりの磁気モーメントは、はかると約 $2.2\mu_{\mathrm B}$ です。原子の向きが全部そろったときの磁化 $M$（1 m$^3$ あたりの磁気モーメント：電-7）と、$\mu_0M$ を求めましょう。鉄の密度は \qty{7.87}{g/cm^3}、モル質量は \qty{55.85}{g/mol}、$\mu_{\mathrm B}=\qty{9.27e-24}{J/T}$ とします。
\phbamen
電-7 で「鉄はおよそ 2 T で飽和する」とした値を、原子の磁石から見積もる場面です。
\phmushi
原子の磁石は全部同じ向きにそろっている（飽和）とします。
\phzu
1 m$^3$ あたりの原子の数 $n_{\text{原子}}=\rho N_{\mathrm A}/M_{\text{モル}}$（$M_{\text{モル}}$ はモル質量）。\\
磁化 $M=n_{\text{原子}}\times2.2\mu_{\mathrm B}$、$\mu_0=4\pi\times10^{-7}$ N/A$^2$。
\phriyu
磁化は、1 m$^3$ の中の小さな磁石の磁気モーメントの合計です（電-7）。
\phtateru
\[ M=\frac{\rho N_{\mathrm A}}{M_{\text{モル}}}\times2.2\mu_{\mathrm B} \]
\phtoku
\[ n_{\text{原子}}=\frac{7.87\times10^6}{55.85}\times6.02\times10^{23}\approx8.48\times10^{28}\ \text{個/m}^3 \]
\[ M=8.48\times10^{28}\times2.2\times9.27\times10^{-24}\approx\qty{1.7e6}{A/m},\qquad \mu_0M\approx\qty{2.2}{T} \]
\phtashikame
単位：(1/m$^3$)×(J/T)＝J/(T·m$^3$)＝A/m（J/T＝A·m$^2$）。$\mu_0M$ は T。\\
大きさ：電-7 の鉄の飽和（約 2 T）と合います。\\
極端な場合：⑧の原子1個の絵（向きのそろった電子4個）なら $4\mu_{\mathrm B}$ で、$\mu_0M\approx\qty{4.0}{T}$ になるはずです。実際が約半分なのは、金属の中では $3d$ の電子もとなりの原子と席を分け合い、帯（原-8）になっているからです。
\kotae $M\approx\qty{1.7e6}{A/m}$、$\mu_0M\approx\qty{2.2}{T}$
\end{phreidai}
\begin{jibun}
例題9 で、原子1個あたりの磁気モーメントがちょうど $1\mu_{\mathrm B}$ だとすると、$\mu_0M$ はいくらになりますか。
\bkjans{$\mu_0\times n_{\text{原子}}\times\mu_{\mathrm B}$ で、約 \qty{0.99}{T}}
\end{jibun}

\label{ph:原-7:ex10}
\begin{phreidai}[水素分子の結合のエネルギー]
水素分子 H$_2$ を2つの水素原子に分けるには、分子1個あたり \qty{4.48}{eV}が要ります。1 mol あたりでは何 kJ ですか。また、室温（300 K）の $k_{\mathrm B}T$ の何倍ですか。
\phbamen
化学の「結合エネルギー」の表の数と、原子1個の世界の eV をつなぐ場面です。
\phmushi
分子の振動の、いちばん低い状態のエネルギー（零点エネルギー：原-6 ⑧）をふくめた、0 K での値とします。
\phzu
$\qty{1}{eV}=\qty{1.60e-19}{J}$、$N_{\mathrm A}=\num{6.02e23}\,\mathrm{/mol}$、$k_{\mathrm B}T=\qty{0.0259}{eV}$。
\phriyu
1 mol は $N_{\mathrm A}$ 個なので、1個のエネルギーに $N_{\mathrm A}$ をかけます。
\phtateru
\[ E_{\mathrm{mol}}=4.48\times1.60\times10^{-19}\times6.02\times10^{23},\qquad \frac{4.48}{0.0259} \]
\phtoku
\[ E_{\mathrm{mol}}\approx\qty{4.32e5}{J/mol}=\qty{432}{kJ/mol},\qquad \frac{4.48}{0.0259}\approx173 \]
\phtashikame
単位：eV×(J/eV)×(1/mol)＝J/mol。\\
大きさ：化学の表の H$_2$ の結合エネルギー（25 ℃での値。約 436 kJ/mol）とほぼ同じです。少し違うのは、温度と、エネルギーの数え方の違いです。$1\,\mathrm{eV}$ が1 mol で約 \qty{96.5}{kJ} にあたることを覚えておくと便利です。\\
極端な場合：結合のエネルギーが $k_{\mathrm B}T$ の170倍もあるので、室温の熱の動きでは、分子はまずこわれません。
\kotae 約 \qty{432}{kJ/mol}。$k_{\mathrm B}T$ の約170倍
\end{phreidai}
\begin{jibun}
水素原子 1 mol をすべて電離する（電子を取り去る）のに要るエネルギーは何 kJ ですか。電離エネルギーは \qty{13.6}{eV}です。
\bkjans{13.6 eV に $N_{\mathrm A}$ をかけて J にすると、約 \qty{1.31e3}{kJ}（約 1310 kJ）}
\end{jibun}

\label{ph:原-7:ex11}
\begin{phreidai}[一酸化炭素の回転の電波（3けたで書く）]
一酸化炭素 CO の分子の回転の段 $\hbar^2/(2I)$ と、回転が $J=0$ から $J=1$ に移るときに吸う電波の周波数を求めましょう。電波望遠鏡の観測と比べるので、3けたで書きます。原子の間の距離は \qty{112.8}{pm}、換算質量は \qty{6.856}{u}、$\qty{1}{u}=\qty{1.661e-27}{kg}$、$\hbar=\qty{1.055e-34}{J.s}$、$h=\qty{6.626e-34}{J.s}$ とします。
\phbamen
宇宙のガスの雲から届く CO の電波の周波数を、分子の形から求める場面です（⑪）。
\phmushi
分子は伸び縮みしない棒（剛体）とし、回ると少し伸びる効果（遠心力）は無視します。
\phzu
$\mu=6.856\times1.661\times10^{-27}$ kg、$r=\qty{112.8e-12}{m}$、$I=\mu r^2$、$E_J=\dfrac{\hbar^2}{2I}J(J+1)$。
\phriyu
回転のエネルギーの段の式（⑪）で、$E_1-E_0=2\times\dfrac{\hbar^2}{2I}=\dfrac{\hbar^2}{I}$。光子のエネルギーが $hf$ です。
\phtateru
\[ I=\mu r^2,\qquad f=\frac{E_1-E_0}{h}=\frac{\hbar^2}{hI} \]
\phtoku
\[ \mu=6.856\times1.661\times10^{-27}\approx\qty{1.139e-26}{kg} \]
\[ I=1.139\times10^{-26}\times(112.8\times10^{-12})^2\approx\qty{1.449e-46}{kg.m^2} \]
\[ \frac{\hbar^2}{2I}=\frac{(1.055\times10^{-34})^2}{2\times1.449\times10^{-46}}\approx\qty{3.84e-23}{J}\approx\qty{2.40e-4}{eV} \]
\[ f=\frac{2\times3.84\times10^{-23}}{6.626\times10^{-34}}\approx\qty{1.16e11}{Hz}=\qty{116}{GHz} \]
\phtashikame
単位：(J·s)$^2$÷(kg·m$^2$)＝J、J÷(J·s)＝Hz。\\
大きさ：電波望遠鏡で観測される CO の線は \qty{115.27}{GHz}（波長 約\qty{2.6}{mm}）で、0.5\,\%ほどの違いで合います。分子はいつも少し振動していて、平均の距離が 112.8 pm より少し長いためです。\\
極端な場合：同じ形で原子が軽い分子（H$_2$ など）ほど $I$ が小さく、段の間かくが大きくなります。
\kotae $\hbar^2/(2I)\approx\qty{2.40e-4}{eV}$、周波数 約 \qty{116}{GHz}
\end{phreidai}
\begin{jibun}
例題11 の CO が、$J=1$ から $J=2$ に移るときに吸う電波の周波数を求めましょう。
\bkjans{$E_2-E_1=(6-2)\times\frac{\hbar^2}{2I}=2\times\frac{\hbar^2}{I}$。$J=0\to1$ の2倍で、約 \qty{232}{GHz}}
\end{jibun}

\label{ph:原-7:ex12}
\begin{phreidai}[原子の小さな磁石を数える]
鉄（$[\mathrm{Ar}]\,3d^6\,4s^2$）、銅（$[\mathrm{Ar}]\,3d^{10}\,4s^1$）、ネオジム（$[\mathrm{Xe}]\,4f^4\,6s^2$）の、満席でない内側の殻（$3d$・$4f$）の、相手のいない電子の数を、フントの規則で数えましょう。
\phbamen
磁石の材料に使われる元素が、原子のどこに磁石のもとをもっているかを調べる場面です。
\phmushi
満席の殻と、外の $s$ の電子（金属の中では動き回る）は数えません。
\phzu
$3d$ は箱5つ（10席）、$4f$ は箱7つ（14席）。
\phriyu
フントの規則（⑧）：箱に1個ずつ同じ向きで入れ、残りを反対の向きで入れます。
\phtateru
相手のいない電子の数は、電子の数と（箱の数×2－電子の数）の小さいほう。
\phtoku
鉄 $3d^6$：上向き5個と下向き1個で、相手のいない電子は $5-1=4$ 個。\\
銅 $3d^{10}$：満席で0個。ネオジム $4f^4$：箱7つに上向き4個で、4個。
\phtashikame
単位：個数。極端な場合：$3d^5$・$4f^7$ で、いちばん多くなります。\\
大きさ：鉄とネオジムは、原子1個が強い磁石のもと。ネオジム磁石（ネオジム・鉄・ホウ素の化合物）は、この2つを組み合わせた材料です。
\kotae 鉄 4個、銅 0個、ネオジム 4個
\end{phreidai}
\begin{jibun}
マンガン（$[\mathrm{Ar}]\,3d^5\,4s^2$）と亜鉛（$[\mathrm{Ar}]\,3d^{10}\,4s^2$）の $3d$ の、相手のいない電子の数を求めましょう。
\bkjans{マンガン 5個（$3d$ の箱5つに1個ずつ）、亜鉛 0個（満席）}
\end{jibun}

%% ======================== まちがえやすいところ ========================
\Needspace{11\baselineskip}\bksection{tsumazuki}{まちがえやすいところをチェック}
\begin{tsumazuki}{電子は、殻の円の上を回っている？}
\bkng{K 殻の電子は、半径 \qty{53}{pm}の円の上を、決まった速さで回っている。}
\bkok{電子は、ふちのぼやけた雲のように広がっている。\qty{53}{pm}は、いちばん見つかりやすい距離。}
殻の絵は、席の数を数えるための絵です。電子の居場所は確率で決まります（②・③）。
\end{tsumazuki}

\begin{tsumazuki}{第3周期の行の長さは18？}
\bkng{M 殻（$n=3$）の席は18個なので、第3周期には18の元素が並ぶ。}
\bkok{第3周期は $3s$ と $3p$ の8個。$3d$ の10個は、$4s$ より高いので、第4周期に回る。}
$2n^2$ は殻の席の数で、周期の長さではありません（⑤の図\ref{ph:原-7:fig-order}）。
\end{tsumazuki}

\begin{tsumazuki}{イオンになるときも、3d から電子が抜ける？}
\bkng{鉄は $4s$ が先にうまったので、Fe$^{2+}$ になるときは、あとから入った $3d$ から2個抜ける。}
\bkok{Fe$^{2+}$ は $[\mathrm{Ar}]\,3d^6$。外側にいる $4s$ の電子から抜ける。}
「うまる順番」と「いちばん外にいる電子」は同じではありません。電子が減ると、$3d$ のほうが低くなります。
\end{tsumazuki}

\begin{tsumazuki}{スピンは、電子が自分で回っていること？}
\bkng{スピンが上向きの電子は、左回りに自転している。}
\bkok{スピンは、電子が生まれつきもつ、角運動量と磁石の性質。こまのように回っているのではない。}
回っていると考えると、電子の表面が光の速さをこえてしまいます（⑥）。
\end{tsumazuki}

%% ======================== 工学ではこう使う ========================
\Needspace{10\baselineskip}\bksection{phgear}{工学ではこう使う}
\begin{phkougaku}[電子の席とスピンを使う技術]
\textbf{医療}：MRI は、体の中の水の陽子のスピンの共鳴（⑥）で、体の断面をうつします。放射線を使わずに、脳や関節の中のようすがわかります。

\begin{phketa}[原子と分子のエネルギーの目安]
\phketaline{いちばん外の電子の電離}{4〜25 eV（アルカリ金属 4〜5、貴ガス 12〜25）}
\phketaline{化学結合}{数 eV（H$_2$ で 4.5 eV、1 eV は1 mol で約 96 kJ）}
\phketaline{分子の回転の段}{$10^{-4}$〜$10^{-3}$ eV（電波）。室温の $k_{\mathrm B}T$ は 0.026 eV}
\phketaline{スピンの共鳴}{電子 28.0 GHz/T、陽子 42.58 MHz/T}
\phboxhead{答えが桁で合っているか}
原子の電子のエネルギーは「13.6 eV に $Z_{\mathrm{eff}}^2/n^2$」の大きさです。いちばん外の電子で何千 eV のような答えが出たら、遮へいを忘れて $Z$ をそのまま入れていないか、確かめましょう。何千 eV は、重い原子のいちばん内側の電子（X線：原-4 ⑩）の大きさです。
\end{phketa}

\textbf{材料}：ESR は、相手のいない電子を見つけて数え、電池や半導体の材料の欠陥を調べます。新しい材料は、まず第一原理計算で電子の席とエネルギーを計算してから作ることが増えています。

\textbf{磁石}：モーターのネオジム磁石は、鉄の $3d$ とネオジムの $4f$ の、向きのそろった電子（⑧）を組み合わせて作ります。

\textbf{光と気体}：ネオンサイン、蛍光灯（水銀とアルゴン）、溶接のガス（アルゴン）は、貴ガスが結びつかない性質と、決まった色の光を出す性質を使います。


\end{phkougaku}
\phstation{材料工学,量子化学,磁性材料,医用画像}{電子配置、量子数、スピン、交換相互作用、電離エネルギー、核磁気共鳴、電子スピン共鳴、第一原理計算}

%% ======================== 章末問題 ========================
\bkendproblems
\bkproblemgroup{基本}{3}
\bkproblem マグネシウム（$Z=12$）・塩素（17）・カリウム（19）の電子配置を書き、K・L・M・N 殻に何個ずつ入るか答えましょう。
\bkpans{Mg：$1s^2\,2s^2\,2p^6\,3s^2$（K2・L8・M2）。Cl：$1s^2\,2s^2\,2p^6\,3s^2\,3p^5$（K2・L8・M7）。K：$1s^2\,2s^2\,2p^6\,3s^2\,3p^6\,4s^1$（K2・L8・M8・N1）}
\bkproblem （1）$n=5$ の殻の席の数を求めましょう。（2）$l=4$（$g$ とよぶ）の席は、1つの $n$ の中に何個ありますか。
\bkpans{（1）$2\times5^2=50$。（2）$2(2\times4+1)=18$}
\bkproblem 塩素とカルシウム（$Z=20$）は、それぞれ電子を何個受け取るか失うかして、どの貴ガスと同じ電子配置のイオンになりやすいですか。イオンの式も書きましょう。
\bkpans{塩素は1個受け取り Cl$^-$（アルゴンと同じ）。カルシウムは $4s$ の2個を失い Ca$^{2+}$（アルゴンと同じ）}
\bkproblem ルビジウム（外の電子は $5s^1$、電離エネルギー \qty{4.18}{eV}）の有効核電荷を求めましょう。
\bkpans{$Z_{\mathrm{eff}}=5\sqrt{4.18/13.6}\approx2.8$}
\bkproblem 7.0 T の研究用の MRI で、陽子が共鳴する周波数と、同じ磁場で電子が共鳴する ESR の周波数を求めましょう（陽子 42.58 MHz/T、電子 28.0 GHz/T）。
\bkpans{陽子 約 \qty{298}{MHz}、電子 約 \qty{196}{GHz}}
\bkproblem ナトリウム 1 mol の原子を、すべて Na$^+$ にするのに要るエネルギーは何 kJ ですか。電離エネルギーは \qty{5.14}{eV}、$\qty{1}{eV}=\qty{1.60e-19}{J}$、$N_{\mathrm A}=\num{6.02e23}\,\mathrm{/mol}$ とします。
\bkpans{5.14 eV に $N_{\mathrm A}$ をかけて J にすると、約 \qty{495}{kJ}}
\bkproblem \qty{1.0}{T}の磁場の中の電子のスピンで、上の向きの数が下の向きの半分になる温度を求めましょう（$\Delta E=\qty{1.16e-4}{eV}$、$k_{\mathrm B}=\qty{8.62e-5}{eV/K}$）。
\bkpans{$e^{-\Delta E/(k_{\mathrm B}T)}=\tfrac12$ より $T=\Delta E/(k_{\mathrm B}\ln2)$。約 \qty{1.9}{K}}
\bkproblemgroup{チャレンジ}{5}
\bkproblem 水素原子の $1s$ の電子が、ちょうど半分の確率で見つかる球の半径 $r_{1/2}$ を求めましょう（例題2 の $F(x)$ を使い、$x=2r_{1/2}/a_0$ を表計算か二分法でさがす）。
\bkpans{$1-e^{-x}\bigl(1+x+\tfrac{x^2}{2}\bigr)=\tfrac12$。$x=2.67$ で 0.499、$x=2.68$ で 0.501 なので $x\approx2.67$。$r_{1/2}\approx1.34a_0\approx\qty{71}{pm}$（いちばん見つかりやすい $a_0$ より外、平均 $1.5a_0$ より内）}
\bkproblem 水素原子の $2s$ の $\psi=\bigl(2-\tfrac{r}{a_0}\bigr)e^{-r/(2a_0)}$ が方程式をみたすことを確かめ、エネルギーを求めましょう（$x=r/a_0$、\qty{27.2}{eV} を単位にすると $-\tfrac12\nabla^2\psi-\tfrac1x\psi=E\psi$）。
\bkpans{$\psi'=e^{-x/2}(\tfrac x2-2)$、$\psi''=e^{-x/2}(\tfrac32-\tfrac x4)$。$\nabla^2\psi=\psi''+\tfrac2x\psi'=e^{-x/2}(\tfrac52-\tfrac x4-\tfrac4x)$。$-\tfrac12\nabla^2\psi-\tfrac1x\psi=e^{-x/2}(-\tfrac14+\tfrac x8)=-\tfrac18\psi$。$1/x$ の項が消えて $E=-\tfrac18\times27.2\approx\qty{-3.40}{eV}$（$-13.6/2^2$ と同じ）}
\bkproblem $^{13}$C$^{16}$O の $J=0\to1$ の電波の周波数を求めましょう。$^{12}$C$^{16}$O は \qty{115.27}{GHz}、距離は同じ、質量は $^{12}$C 12.00 u、$^{13}$C 13.00 u、$^{16}$O 15.99 u とします。
\bkpans{周波数は換算質量に反比例する。換算質量は $^{12}$C$^{16}$O で約 6.855 u、$^{13}$C$^{16}$O で約 7.170 u。$115.27\times6.855\div7.170\approx\qty{110.2}{GHz}$}

\begin{furikaeri}
\begin{checklist}
\item 周期表の横の並びの数（2・8・8・18…）が、原子の中の電子の「席」の数で決まることを説明できる。\bkfurref{ph:原-7:sec1}
\item 量子数で席を数えて $2n^2$ を導き、ナトリウム・鉄・銅などの電子配置を書ける。\bkfurref{ph:原-7:sec4}
\item 水素原子の電子の雲から、電子が見つかりやすい距離や見つかる確率を求め、電離エネルギーから有効核電荷を見積もれる。\bkfurref{ph:原-7:sec3}
\item 磁場の中の電子や陽子の向きが2つに分かれることから、ESR や MRI の周波数を求められる。\bkfurref{ph:原-7:sec6}
\item 3次元のシュレディンガー方程式に水素原子の $1s$ の波動関数を代入して $-13.6$ eV を確かめ、化学結合や分子の回転のしくみを読める。\bkfurref{ph:原-7:sec10}
\end{checklist}
\end{furikaeri}
\bkfukidashi{mina}{ネオンが結びつかないのは、電子の席がちょうど満席だったからなんだね。2・8・8 の数も、MRI も、磁石も、「1つの席に1個」の決まりから出てくるなんて、びっくり！}
